\documentclass[11pt]{article}
\usepackage[T1]{fontenc}
\usepackage{lmodern,amsmath,amssymb,amsthm,mathtools}
\usepackage[margin=0.9in]{geometry}
\usepackage{booktabs,longtable,array,microtype,xcolor,fancyhdr}
\usepackage{xurl}
\usepackage[colorlinks=true,linkcolor=blue!45!black,urlcolor=blue!45!black,citecolor=blue!45!black]{hyperref}
\usepackage{enumitem,needspace}
\setlist{itemsep=3pt,topsep=5pt}
\allowdisplaybreaks
\setlength{\emergencystretch}{3em}
\pagestyle{fancy}\fancyhf{}
\fancyhead[L]{\small Con(ZFC): oriented power certificates}
\fancyhead[R]{\small 25 October 2026}
\fancyfoot[C]{\thepage}
\setlength{\headheight}{14pt}
\newtheorem{theorem}{Theorem}[section]
\newtheorem{lemma}[theorem]{Lemma}
\newtheorem{corollary}[theorem]{Corollary}
\newtheorem{proposition}[theorem]{Proposition}
\theoremstyle{definition}\newtheorem{definition}[theorem]{Definition}
\theoremstyle{remark}\newtheorem{remark}[theorem]{Remark}
\newcommand{\Z}{\mathbb Z}
\newcommand{\Q}{\mathbb Q}
\newcommand{\R}{\mathbb R}
\newcommand{\N}{\mathbb N}
\newcommand{\Con}{\operatorname{Con}}
\newcommand{\sgn}{\operatorname{sgn}}
\newcommand{\height}{\operatorname{ht}}
\newcommand{\src}[1]{\texttt{\detokenize{#1}}}
\title{Oriented Power Pooling and Polynomial Root Centering\\for Diophantine Consistency Statements}
\author{Research prepared for Vladimir Reshetnikov and ProveIt\\
\small AI-assisted development, with separately assigned mathematical and algebraic reviews}
\date{25 October 2026}
\hypersetup{pdftitle={Oriented Power Pooling and Polynomial Root Centering for Diophantine Consistency Statements},pdfauthor={AI-assisted research prepared for Vladimir Reshetnikov},pdfsubject={Explicit Diophantine equations equivalent to Con(ZFC)}}

\begin{document}
\maketitle
\begin{abstract}
We give a complete construction that combines all seventeen power requests
in ProveIt's direct C222 arithmetic encoding into one rational-integrality
certificate. Multiplication by the base orients the error terms even when
bases have different signs. Zero bases are allowed in the local theorem;
the retained arithmetic equations subsequently exclude them. This avoids
a circular assumption about the radix and removes nine certificate
coordinates. We then apply explicitly listed acyclic unit-pivot
substitutions. The resulting fixed integer polynomials include forms in
87 unrestricted integers of degree 8,388,612, in 88 of degree 65,540,
in 89 of degree 32,768, and
moderate-degree alternatives listed below. Their arithmetic circuits use
only integer literals of absolute value at most 256; their expanded
coefficients are larger and are bounded separately.

A separate compact result resolves the earlier square-transport question:
an inverse Lucas correction permits the transport weight $b^2$, followed
by an integral polynomial root center. It produces a ten-positive-variable
polynomial of exact degree 224 in 90 binary operations, together with
85-operation degree-232 companions. This branch retains large calibrated
coefficients and has a different domain from the direct integer branch.

For every delivered equation, composition with the specifically identified
inherited refutation calibration establishes in ZFC that absence of a
solution in its stated domain is equivalent to $\Con(\mathrm{ZFC})$.
The new unbounded arguments are ordinary mathematical proofs. Independent
finite checks authenticate the formulas, witness transformations, degrees,
and circuits; no new end-to-end proof-assistant compilation is claimed.
\end{abstract}

\tableofcontents

\section{Results, comparison, and the meaning of simplicity}

The task is to specify an actual polynomial with fixed integer
coefficients, rather than an unspecified instance of a universal
representation theorem. The delivered direct family has the following
precise property. For every listed pair $(n,d)$ there is a completely
specified $P_{n,d}\in\Z[X_0,\ldots,X_{n-1}]$, of exact total degree $d$,
such that
\begin{equation}
\mathrm{ZFC}\ \vdash\quad
\bigl(\neg\exists\mathbf X\in\Z^n\ P_{n,d}(\mathbf X)=0\bigr)
\ \longleftrightarrow\ \Con(\mathrm{ZFC}).
\label{eq:target}
\end{equation}
Here consistency means syntactic consistency for the effective ZFC
presentation used in the retained calibration. None of the construction's
integer coefficients depends on deciding whether that theory is consistent.

\subsection{The direct family}

\begin{center}\small
\begin{tabular}{rrrrr}
\toprule
Integer variables & Exact degree & Residuals & Operations & $H$\\
\midrule
137 & 12 & 74 & 2,548 & 55\\
121 & 16 & 58 & 2,502 & 84\\
118 & 22 & 55 & 2,523 & 58\\
111 & 32 & 48 & 2,502 & 107\\
106 & 38 & 43 & 2,457 & 147\\
104 & 64 & 41 & 2,451 & 223\\
99 & 128 & 36 & 2,483 & 556\\
94 & 256 & 31 & 2,468 & 788\\
93 & 328 & 30 & 2,448 & 788\\
89 & 32,768 & 26 & 2,453 & 82,436\\
88 & 65,540 & 25 & 2,450 & 164,875\\
87 & 8,388,612 & 24 & 2,447 & 21,104,473\\
\bottomrule
\end{tabular}
\end{center}


In the table, $H$ is an integer such that every collected coefficient has
absolute value strictly below $2^H$. These are proved upper bounds, not
exact coefficient heights. Operations count binary additions, subtractions,
and multiplications in the supplied circuit, with a square counted as a
multiplication and a multiplication by a constant also counted. They are
not bit-complexity estimates. All listed variables are independently
quantified over $\Z$, and every variable occurs in the actual polynomial;
the independent checker certifies the latter with nonzero modular partial
derivatives, beyond mere occurrence in the computation graph.

The 87-variable form is the smallest-dimensional delivered direct form;
its degree is substantially higher than the other profiles. It is obtained
from the 88-variable form by one additional unit pivot. The 88-variable
form has degree 65,540, and the 89-variable form reduces the degree to
32,768 at the cost of one coordinate.
The 111-variable degree-32 form is a useful moderate-degree choice, with
a comparatively small proved coefficient-height bound. A reader whose
priority is short evaluation rather than unrestricted integer variables
should instead examine the compact positive-domain companion in
Section~\ref{sec:compact-center}.

\subsection{What is improved, and what is inherited}

The comparison snapshot is the actual repository commit
\begin{center}\small\texttt{0c48d332df6b0e50d4e5b6251c531898f5682c72}.\end{center}
This matters because the intake contained newer constructions than the
frontier summarized in the older top-level README. The immediate direct
parent is package 158, with 225 integer variables and 162 quadratic
residuals. Package 168 already supplied a two-base reciprocal certificate
and a 215-variable degree-ten companion. Its final discussion suggested
oriented weights for more bases. We credit that suggestion. The present
contribution is the completed five-base theorem, including temporary zero
bases, the simultaneous bootstrap, the actual polynomials, and their
composition with certified coordinate elimination.

Package 152 had already applied acyclic monic elimination to a different
direct parent. Its relevant reported points were
\[
(144,12),\ (125,16),\ (117,32),\ (111,64),\ (102,256),\ (94,65540),
\]
where each pair is (unrestricted integer variables, exact degree).
Our family improves the dimension at these selected degrees or provides
additional tradeoffs; it does not improve every metric simultaneously.
For example the supplied factored evaluators can require more binary
operations than package 152's evaluators. The comparison is to authenticated
repository constructions, not a proof of a global optimum.

The direct route has no native coefficient tuple and does not use a K215
machine simulation, Jones's universal polynomial, or Sun's integer-domain
reduction. Its starting point is the repository's numerical proof calculus
C222. The compact ten-positive-variable companion retains the separate
native/K215 calibration and four large fixed coefficients. These two
ancestries and domains are kept explicit throughout.

The existence of Diophantine descriptions for recursively enumerable
relations is classical \cite{matiyasevich,matbook}. That general theorem
does not identify the best explicit equation for this particular
consistency statement. Nor does a universal variable bound, such as the
recent ten-integer-variable result \cite{sun}, settle the simultaneous
optimization of degree, coefficient size, circuit size, and calibration
complexity.

\section{The frozen parent and the exact finite definition}
\label{sec:definition}

Let $R_0,\ldots,R_{161}\in\Z[x_0,\ldots,x_{224}]$ be the literal
quadratic residuals in Appendix~\ref{app:parent}. The parent polynomial is
\[
P_{225}=\sum_{r=0}^{161}R_r^2.
\]
\begin{samepage}
Its byte-identical coefficient file has SHA-256
\begin{center}\scriptsize\texttt{fee84c408828aba3e7ceebe9a37f58509327063aa7fd078fd49c80adf3af4e56}.\end{center}
The original Git blob is
\texttt{425a70a4960e0c60ebedfcae641f08be15dbdcc0}.
The parent is not a parameter to be selected later: its complete finite
coefficient list is frozen and included.
\end{samepage}

All newly pooled systems are specified by the same following finite recipe.
Delete the ten old coordinates
\begin{equation}
149,150,151,152,153,154,221,222,223,224
\label{eq:deleted-variables}
\end{equation}
and the ten old rows
\begin{equation}
105,106,107,108,109,110,158,159,160,161.
\label{eq:deleted-rows}
\end{equation}
Retain the other coordinates in increasing old-index order, relabeling them
$u_0,\ldots,u_{214}$. Retain the other residuals in increasing old-row
order. Append the one, two, or three new integer coordinates and rows
specified in Section~\ref{sec:groupings}. This gives the raw all-five,
staged, and paired systems in respectively 216, 217, and 218 variables.

The difference coordinates $D_i=X_i-AY_i$ are already present in the
parent. The entire finite request table is as follows. The symbols $X_i$
are reconstructed as $D_i+Z_i$, where the retained multiplication rows
give $Z_i=AY_i$. Set $A=x_{142}+1$, $q=x_{102}-1$, $t=x_0$.

\begin{center}\small
\begin{tabular}{rrlllll}
\toprule
Request & Group & Index $k_i$ & Output $m_i$ & $D_i$ & $Y_i$ & $Z_i$\\
\midrule
0 & 0 & $x_0$ & $x_5$ & $x_{146}$ & $x_{147}$ & $x_{148}$\\
1 & 1 & $x_1$ & $x_{18}$ & $x_{182}$ & $x_{183}$ & $x_{184}+1$\\
2 & 2 & $x_1$ & $x_{102}-1$ & $x_{182}$ & $x_{183}$ & $x_{184}+1$\\
3 & 4 & $x_0+x_3$ & $x_{103}$ & $x_{155}$ & $x_{156}$ & $x_{157}$\\
4 & 3 & $x_0$ & $x_{104}$ & $x_{146}$ & $x_{147}$ & $x_{148}$\\
5 & 4 & $x_9-1$ & $x_{108}$ & $x_{158}$ & $x_{159}-1$ & $x_{160}$\\
7 & 4 & $x_{10}$ & $x_{111}$ & $x_{161}$ & $x_{162}$ & $x_{163}$\\
8 & 4 & $x_{11}$ & $x_{114}$ & $x_{164}$ & $x_{165}$ & $x_{166}$\\
9 & 4 & $x_{12}$ & $x_{117}$ & $x_{167}$ & $x_{168}$ & $x_{169}$\\
10 & 4 & $x_{13}$ & $x_{120}$ & $x_{170}$ & $x_{171}$ & $x_{172}$\\
11 & 3 & $x_5$ & $x_{121}$ & $x_{173}$ & $x_{174}$ & $x_{175}$\\
12 & 4 & $x_{14}$ & $x_{124}$ & $x_{176}$ & $x_{177}$ & $x_{178}$\\
13 & 4 & $x_{15}$ & $x_{127}$ & $x_{179}$ & $x_{180}$ & $x_{181}$\\
14 & 3 & $x_1$ & $x_{128}$ & $x_{182}$ & $x_{183}$ & $x_{184}+1$\\
15 & 4 & $x_{16}$ & $x_{132}$ & $x_{185}$ & $x_{186}$ & $x_{187}$\\
16 & 4 & $x_{17}$ & $x_{135}$ & $x_{188}$ & $x_{189}$ & $x_{190}$\\
17 & 4 & $x_{100}$ & $x_{138}$ & $x_{191}$ & $x_{192}$ & $x_{193}$\\
\bottomrule
\end{tabular}
\end{center}


For each of the five base groups, define the following polynomial
abbreviations, which introduce no new unknowns:
\begin{align}
T_j&=2An_j-n_j^2-1,\nonumber\\
H_j&=\sum_{i\in I_j}Y_i(m_i-D_i-Z_i),\nonumber\\
S_j&=\sum_{i\in I_j}\bigl((m_i-D_i)^2-Y_i^2\bigr),\nonumber\\
E_j&=S_j-2n_jH_j.\label{eq:group-definitions}
\end{align}
The five bases have request multiplicities $1,1,1,3,11$, respectively:
\[
n_0=x_3+1=F^2+1,\quad n_1=x_4=B,\quad n_2=x_{194}=BH,
\quad n_3=q,\quad n_4=q+1,
\]

Finally apply the finite pivot lists in Appendix~\ref{app:pivots}. In each
list the row and variable numbers refer to the indicated raw pooled
system, not to the old parent. A pair $(r,v)$ replaces the coordinate
$u_v$ by the integer polynomial obtained by solving its unit-linear row
$R_r=0$. Definitions are recursively substituted in their acyclic
dependency order. Keep the remaining coordinates in increasing raw-index
order, substitute into every unremoved residual, and sum their squares.
This is a finite, complete definition of every $P_{n,d}$ in the results
table, independently of running the provided software.

The JSON circuits are a more convenient equivalent presentation. Their
internal registers are polynomial abbreviations; only their declared input
variables are quantified. They contain integer constants and the operations
$+,-,\times$, with no divisions, factorials, variable exponents, or
unevaluated universal-polynomial parameters.

\section{Oriented pooling with temporarily zero bases}
\label{sec:oriented}

\subsection{Canonical signed Pell coordinates}

For an integer $A$ with $|A|>1$, define the canonical Pell coordinates by
$(X_0,Y_0)=(1,0)$ and
\[
X_{k+1}=AX_k+(A^2-1)Y_k,\qquad
Y_{k+1}=X_k+AY_k
\quad(k\ge0),
\]
and $X_{-k}=X_k$, $Y_{-k}=-Y_k$. These definitions work for either sign
of $A$. They give $X_k^2-(A^2-1)Y_k^2=1$.
For $T=2An-n^2-1$ one has $(n-A)^2\equiv A^2-1\pmod T$.
For $n\ne0$ and $k\in\Z$, or for arbitrary $n$ and $k\ge0$,
the Pell group homomorphism therefore gives
\begin{equation}
X_k+(n-A)Y_k\equiv n^k\pmod T.\label{eq:pell-congruence}
\end{equation}
When $n\ne0$, negative exponents in this congruence mean modular
inverses. They exist because $T\equiv-1\pmod n$. No rational power
occurs in the final polynomial.

\begin{theorem}[One certificate for all active bases]
\label{thm:all-bases}
Let $W\ge258$ and $|A|>W^W$ be integers. There are at most five blocks
of requests and at most seventeen requests in total. In block $j$ the
base is an arbitrary integer $n_j$. A request consists of integers
$m_i,k_i$ and the canonical signed pair $(X_{k_i},Y_{k_i})$ at $A$.
Assume
\[
n_j^2<W,\qquad m_i^2<W,\qquad k_i^2<W.
\]
Set $D_i=X_{k_i}-AY_{k_i}$, $Z_i=AY_{k_i}$, and define
$T_j,H_j,S_j,E_j$ by \eqref{eq:group-definitions}. Then every $T_j$ is
nonzero. The polynomial equation
\begin{equation}
\mathcal R=
\sum_j n_jE_j\prod_{\ell\ne j}T_\ell
-z\prod_jT_j=0
\label{eq:oriented}
\end{equation}
has an integer witness $z$ if and only if every request at a nonzero
base is exact: $m_i=n_j^{k_i}$ in $\Q$. When $|n_j|\ge2$, exactness
forces $k_i\ge0$. A zero-base block contributes zero and is deliberately
unconstrained at this stage.
\end{theorem}

\begin{proof}
Put $b=\lfloor\sqrt{W-1}\rfloor$, so $b\ge16$ and every absolute
base, output, and index is at most $b$. If $n_j=0$, then $T_j=-1$.
Otherwise
\begin{equation}
|T_j|\ge2|A||n_j|-n_j^2-1\ge2|A|-W>W^W,
\qquad \sgn\frac{n_j}{T_j}=\sgn A.
\label{eq:orientation}
\end{equation}
In particular no denominator vanishes.

Write $\delta_i=m_i-D_i-n_jY_i$. Direct expansion, using $Z_i=AY_i$,
gives the exact identity
\begin{equation}
\sum_{i\in I_j}\delta_i^2
=E_j-T_j\sum_{i\in I_j}Y_i^2.
\label{eq:error-identity}
\end{equation}
Dividing \eqref{eq:oriented} by the nonzero product of the $T_j$,
and subtracting the integral correction $\sum n_jY_i^2$, shows that
\[
Q_0=\sum_{j}\sum_{i\in I_j}\frac{n_j\delta_i^2}{T_j}
\]
is an integer. Terms from zero-base blocks vanish.

Let $J=\{j:n_j\ne0\}$, $s=|J|\le5$, and set
$C=\prod_{j\in J}n_j^b$, with the empty product equal to 1.
This is a nonzero integer. For an active request define
\[
d_i=Cm_i-\frac{C}{n_j^b}n_j^{b+k_i}.
\]
This is an integer: the quotient is an integer, and $b+k_i\ge0$.
Congruence \eqref{eq:pell-congruence} gives
$C\delta_i\equiv d_i\pmod {T_j}$. Hence
\[
Q_1=\sum_{j\in J}\sum_{i\in I_j}\frac{n_jd_i^2}{T_j}
\]
differs from $C^2Q_0$ by an integer and is itself an integer. No
coprimality between different denominators is used.

If $J$ is empty the conclusion is immediate. Otherwise
\[
|C|\le b^{sb},\qquad
|d_i|\le b^{sb+1}+b^{(s+1)b}\le2b^{6b}.
\]
Writing $r\le17$ for the number of requests gives
\begin{equation}
\sum_{j\in J}\sum_{i\in I_j}|n_j|d_i^2
\le4r b^{12b+1}<W^{6b+2}\le W^W.
\label{eq:bound}
\end{equation}
For the strict inequality use $4r\le68<W$, $b^2<W$, and $b<W$.
For the last inequality use $6b+2\le b^2+1\le W$, valid for $b\ge16$.
Equations \eqref{eq:orientation} and \eqref{eq:bound} imply
$|Q_1|<1$. Thus $Q_1=0$. Every nonzero summand has the same sign,
so every $d_i=0$. Division by $C$ in $\Q$ proves
$m_i=n_j^{k_i}$. If $|n_j|\ge2$ and $k_i=-h<0$, this would imply
$m_in_j^h=1$, impossible for integers with $h>0$. Unit bases are
handled without incorrectly inferring a sign for their indices.

Conversely, exact powers at the nonzero bases give
$T_j\mid\delta_i$ by \eqref{eq:pell-congruence}, and then
$T_j\mid E_j$ by \eqref{eq:error-identity}. Therefore
\[
z=\sum_{j:n_j\ne0}\frac{n_jE_j}{T_j}\in\Z
\]
satisfies \eqref{eq:oriented}. Zero-base terms contribute nothing.
\end{proof}

\begin{remark}[Why the weights and zero cases matter]
Without the factors $n_j$, terms from bases of different signs can
cancel. Merely squaring the denominators would not repair the proof:
a congruence modulo $T_j$ gives no corresponding control modulo $T_j^2$.
The factors $n_j$ orient the terms while making the zero-base case
harmless. The active set, floor, and variable powers above are proof
devices only. Equation~\eqref{eq:oriented} is an ordinary polynomial.
\end{remark}

\section{Recovering the five nonzero bases without circularity}
\label{sec:bootstrap}

The parent supplies the following frame before interpreting the power
requests: $W=2+\sum a_\nu^2\ge258$, with every actual base, index,
and proposed output among the 36 explicitly authenticated affine atoms;
$|A|>W^W$; and the canonical signed Pell pairs. Its orientation-first
argument uses retained frame, norm, residue, and cap rows. It does not
use the ten deleted power-certificate rows. The exact atom identity and
the arithmetic proof dependencies are included in the package.

Use the semantic coordinates
\[
F=x_2,\quad t=x_0,\quad V=x_5,\quad B=x_4,\quad
N=x_{18},\quad q=x_{102}-1,\quad M=x_{128},\quad BH=x_{194}.
\]
The following deductions are valid at every zero of the new system.

First, $R_0=0$ gives $x_3=F^2$, and $R_{139}=0$ gives
$t=16+x_{195}^2$. The common remainder row $R_{154}=0$ has the form
\begin{equation}
M^2=1+4\left(S^2+\sum_{i=1}^{11}R_i'^2\right)
+a^2+b^2+c(c+1).
\label{eq:normalization}
\end{equation}
Every summand on the right is nonnegative for integer arguments. Thus
$M\ne0$ and $|S|<|M|/2$. Rows $R_{155},R_{156}$ say
\[
Bv_B=M,\qquad FM=v_BV+S.
\]
Consequently $B,v_B\ne0$. With the integer $s_0=FB-V$, these equations
give $S=v_Bs_0$ and
\begin{equation}|s_0|<|B|/2.\label{eq:half-width}\end{equation}
None of these conclusions used an interpreted power.

The first base $F^2+1$ is positive. Theorem~\ref{thm:all-bases} therefore
already recovers $V=(F^2+1)^t\ge1$. Put $u=|F|$, $b_0=|B|$.
If $b_0=1$, integrality and \eqref{eq:half-width} give $s_0=0$,
whereas $V=FB\le u<(u^2+1)^t=V$, a contradiction. If $u\ge1$ and
$b_0\le u$, then
\[
V=FB-s_0<(u+\tfrac12)b_0
\le u^2+u/2<(u^2+1)^2\le V,
\]
also a contradiction. Thus
\begin{equation}
|B|\ge2,\qquad |F|<|B|,\qquad V<|B|^2.
\label{eq:bootstrap-bounds}
\end{equation}
The last inequality follows from $u\le b_0-1$ and
$V<(u+1/2)b_0$.

The unchanged row $R_{138}=0$ gives
\[
BH=B\bigl(2(x_9+x_{10}+\cdots+x_{17})-1\bigr).
\]
The factor in parentheses is odd, hence nonzero even before the CRT
representatives have been shown nonnegative. Therefore $|BH|\ge2$.
The theorem now recovers $N=B^x$, $q=(BH)^x$, and $x\ge0$.
If $x=0$, then $q=1$; if $x>0$, then $|q|\ge2$. In both cases
$q\ne0,-1$, so both remaining groups are active. All seventeen exact
powers have now been recovered.

The temporary case $x=0$ is excluded by $M=q^x=1$, contradicting
$Bv_B=M$ and $|B|\ge2$. The retained ratio for $q^t$ then gives
$x\ge t$, so $|q|\ge2^{16}$; the subsequent factorial, signed
remainder, CRT, and Horn-flow interpretation is the parent's unchanged
argument. In particular factoriality is not assumed to justify the
signs needed for its own derivation.

\subsection{Witness maps and the inherited consistency endpoint}

For each deleted singleton, restore its modulus $T_j$ and its old quotient
$v_j=-\delta_j/T_j$. These are integers on the new zero set by the exact
power congruence. For an old same-base pool restore
\[
L_j=E_j/T_j,\qquad K_j=2H_j+(2A-n_j)L_j.
\]
These satisfy its original two rows. The independent finite checker proves
that the ten deleted coordinates are private to the ten deleted rows.
Thus the restoration gives an actual zero of $P_{225}$.

In the forward direction there is even a polynomial witness map. At an
old zero, $\delta_j+T_jv_j=0$ and
\[
E_j/T_j=T_jv_j^2+Y_j^2
\]
for each singleton. The all-five witness can therefore be chosen as
\[
z=\sum_{j=0}^{2}n_j(T_jv_j^2+Y_j^2)
+qL_q+(q+1)L_{q+1}.
\]
The reverse restoration is integral on the solution set but is not a
polynomial automorphism on arbitrary tuples. No bijection of all old
certificate witnesses, or preservation of their multiplicities, is claimed.

The inherited logical interfaces are
\begin{align*}
\src{KEndpoint.c222_iff_zf_refutation}:&\quad
\mathrm{C222}(\bot)\ \longleftrightarrow\ \mathrm{Ref}_{\mathrm{ZF}}(\bot),\\
\src{ConstructibleChoice.zf_refutation_iff_zfc_refutation}:&\quad
\mathrm{Ref}_{\mathrm{ZF}}(\bot)\ \longleftrightarrow\ \mathrm{Ref}_{\mathrm{ZFC}}(\bot).
\end{align*}
The parent arithmetic proof includes its square-shell recoding of the
numerical calculus. The logical calibration uses ordinary first-order
soundness and completeness and set-model constructions, including an
internally defined well-founded part and constructibility. It should not
be described as merely a newly implemented primitive-recursive translator
of refutations. The inherited audit records why these constructions are
formalizable in ZFC and do not assume an externally well-founded model
or the consistency statement being represented. Composition gives
\eqref{eq:target}, without deciding either side.

\section{Grouping the requests to control degree}
\label{sec:groupings}

The one-pool formula \eqref{eq:oriented} has residual degree 12 in the
raw coordinates, hence gives a degree-24 sum of squares in 216 integers.
Two alternative groupings reduce degree before coordinate elimination.

\paragraph{Staged grouping.}
For the first three blocks use
\[
\mathcal R_{012}=\sum_{j=0}^{2}n_jE_j
\prod_{\substack{0\le\ell\le2\\\ell\ne j}}T_\ell
-z_0T_0T_1T_2.
\]
For the two consecutive bases use
\[
\mathcal R_{34}=E_3T_4+E_4T_3-z_1T_3T_4.
\]
The first group uses Theorem~\ref{thm:all-bases}. The bootstrap through
$q=(BH)^x$ then makes $q,q+1$ nonzero and of the same sign, so
$T_3,T_4$ have the same sign. The unweighted two-base version of the
same rational-integrality proof applies, as in package 168. The two
residuals have exact degrees 8 and 5. Keeping the other 152 quadratic
rows gives 154 residuals in 217 integers and exact degree 16.

\paragraph{Paired grouping.}
Keep the first singleton quotient but substitute its modulus directly:
\[
\mathcal R_0=\delta_0+z_0T_0.
\]
Use an oriented pair for $B,BH$ and an unweighted pair for $q,q+1$:
\[
\mathcal R_{12}=n_1E_1T_2+n_2E_2T_1-z_1T_1T_2,
\qquad
\mathcal R_{34}=E_3T_4+E_4T_3-z_2T_3T_4.
\]
The first equation is just the parent's first power row with its modulus
definition substituted. It recovers $V$ at the already positive base
$F^2+1$. The oriented pair then recovers $N$ and $q$, and the last
pair completes the bootstrap. The residual degrees are 3, 6, and 5,
giving 155 rows in 218 integers and exact degree 12. The two first
nontrivial bases share the same indexed Pell pair, which makes their
expanded pair numerator particularly short.

These are three exact presentations of the same existence problem, not
an assertion that a higher-degree profile is preferable for every purpose.
The coefficient files also retain the expanded individual residuals as
an independent definition. We do not claim to have collected every
coefficient of the final high-degree sums of squares.

\section{Eliminating defined coordinates with exact certificates}
\label{sec:projection}

\begin{theorem}[Acyclic unit-pivot elimination]
Let $Q_0,\ldots,Q_{m-1}\in\Z[u_0,\ldots,u_{N-1}]$. Select distinct
rows and distinct coordinates so that each selected row is
$Q_r=\varepsilon u_v+f_v$, where $\varepsilon=\pm1$ and $u_v$ does
not occur in $f_v$. Suppose the directed dependency graph of the
selected definitions is acyclic. Recursively substitute
$u_v=-\varepsilon f_v$ into every unselected row, obtaining
$\widehat Q_r$ in the remaining $n$ coordinates. Then projection to the
remaining coordinates is a bijection from the old simultaneous zero
set to the new simultaneous zero set, with an integer polynomial inverse.
In particular
\[
\exists\mathbf u\in\Z^N\ \sum Q_r(\mathbf u)^2=0
\quad\longleftrightarrow\quad
\exists\mathbf X\in\Z^n\ \sum\widehat Q_r(\mathbf X)^2=0.
\]
\end{theorem}

\begin{proof}
Use a topological order of the selected definitions. Every removed
coordinate is uniquely determined by earlier definitions and the retained
coordinates; all coefficients remain integral because the pivot is a unit.
The selected rows vanish identically after substitution. Every unselected
row becomes precisely its stated substituted polynomial. These facts
prove both inverse identities on the corresponding zero sets. A sum of
integer squares is zero exactly when each summand is zero.
\end{proof}

The selected lists are finite certificates, so reproduction does not
depend on rediscovering them with an optimizer. Some pairs solve a
contracted singleton power row for its difference coordinate or output;
that is legitimate whenever the explicit unit-pivot test holds. There
is no requirement that a row originally called a ``power'' row be
immune to a later purely polynomial elimination.

\subsection{Degree, height, and genuine variable occurrence}

Assign degree 1 to every retained input. Each eliminated definition gets
the maximum weighted degree of its source monomials; each remaining
residual gets the same weighted upper bound. Let $D$ be the maximum.
To prove that this bound is exact, evaluate the homogeneous component of
degree $D$ of a residual after substituting the deterministic input
direction $a_v=v+3$, modulo $p=1,000,000,007$. The supplied certificate
records a nonzero result. An integer polynomial with a nonzero reduction
under one such specialization is not the zero polynomial. Thus a residual
has degree exactly $D$. Over $\R$, the sum of squares of the highest
homogeneous components cannot cancel, so the single polynomial has exact
degree $2D$. The proof does not rely on treating a circuit's syntactic
degree as automatically exact.

For coefficient bounds, recursively assign every input coefficient norm
1. A definition $\sum c_m\prod u_i$ receives the bound
$\sum |c_m|\prod L_i$ for its coefficient $\ell^1$ norm. If $L_r$
bounds the norm of the substituted residual, then
$\sum_rL_r^2$ bounds the norm, hence the height, of the final polynomial.
Its bit length is the table's strict exponent $H$.

Finally, circuit reachability alone cannot exclude cancellation of an
input from the actual polynomial. The independent verifier differentiates
the complete circuit modulo $p$ and records an evaluation at which each
partial derivative is nonzero. Those certificates prove actual dependence
on every declared input. They are algebraic nonvanishing witnesses, not
probabilistic evidence about the existence of zeros.

\section{A compact positive-domain companion}
\label{sec:compact-center}

This companion has ten strictly positive integer variables.  Its domain
must be kept separate from the unrestricted-integer equations in the main
construction.  It combines the unsquared strong trace of the pinned
\texttt{163-trace-radical} contribution with the inverse-Lucas correction
of \texttt{164-unsquared-factors-linear-selectors}~\cite{trace163,inverse164}.  A new square transport
weight then permits a polynomial root center and lowers the exact degree.

\begin{center}
\begin{tabular}{lrrr}
\hline
Circuit & Variables & Degree & Binary operations\\
\hline
Trace parent & 10 & 232 & 87\\
Inverse correction with trace & 10 & 232 & 85\\
Square transport with trace & 10 & 232 & 85\\
Centered companion & 10 & 224 & 90\\
\hline
\end{tabular}
\end{center}

The degree 224 point is a measured degree/cost tradeoff.  The claimed result
is neither a global minimum nor an assertion of literature-wide priority.
This compact branch retains its inherited native/K215 calibration.  Its
ordinary proof is separate from its finite algebraic checker.

\subsection{Entire polynomial}

The independent variables are
\[
J,F,\alpha,i,T,w,s,n,\eta,\zeta\in\mathbb Z_{>0}.
\]
The four fixed coefficients \(\beta,K_0,M_C,M_F\) are the whole
indexed BFS943/radix67 tuple, with \(\beta=2^{590069}-1\).  Their complete
finite recipe, 436-character seed, numerical payloads, and independent
reconstruction program are included.  All raw and compressed hashes agree
with the pinned parent.  The largest port has 592,884 bits; the four raw
payloads occupy 295,388 bytes.  Every symbol below is a polynomial
abbreviation, not another quantified variable:
\begin{gather*}
h=\beta J,\quad q=h+1,\quad C_0=q-F-\alpha,\\
X=qw,\quad Y=qs,\quad k=\eta+\zeta,\quad c=kY+\eta,\\
A=Y(X+1)+1,\quad \Delta=A^2-1,\quad H=4A-5,\\
v=XY^2,\quad \ell=vk=(XY)(kY),\quad E=1+\ell(\ell+k),\\
R=[h(q-F)+\alpha](q^2-1)+(M_C+qM_F)J,\\
B=(K_0+w)C_0-F,\quad U=ic^2,\quad a=\Delta U^2E,\\
F_0=1+a,\quad f=1+2a,\quad G=EF_0,\quad Q=4\Delta aF_0,\\
b=cT-R,\quad V=fb-c,\quad K=X(3c-2X)+2,\\
e_6=6(k-R-1)-\ell(k^2-4),\quad D=6q^6,\quad M=hDH.
\end{gather*}
For a weight \(L\), set
\[
Z_0=hD\bigl[H(n-LB)+K\bigr]-H\bigl[DB+h(e_6-6q^4s)\bigr].
\]
The centered companion uses
\[
L=b^2,\qquad C=AUE,\qquad S=CV,\qquad Z=Z_0-MS
\]
and the single equation
\begin{equation}
\boxed{\ P=Q\bigl(Z^2-M^2GV^2\bigr)+(M^2G-Z^2)=0.\ }
\label{eq:compact-centered}
\end{equation}
The two 85-operation controls use \(Z=Z_0\) and respectively
\(L=bV\) or \(L=b^2\).  All three literal DAGs are supplied.  Fixed
powers are multiplication gates; no variable exponent, recurrence, floor,
or unspecified coefficient occurs in the polynomial.

For proof notation only, put
\[
e_*=\frac{e_6}{6},\quad \rho=\frac{e_*}{q^6},\quad
\sigma=\frac{s}{q^2},\quad t=\frac{B}{h},\quad W=1+hL,
\quad\kappa=\frac{K}{H},\quad z=\frac{Z}{M}.
\]
Without yet asserting integrality, the centered carrier identity is
\begin{equation}
z=n-S+\kappa+\sigma-Wt-\rho.
\label{eq:compact-carrier}
\end{equation}
For an uncentered control, omit \(S\).

\subsection{The square transport theorem}

The root extraction, rational adjacent-Pell descent, signed folding, and
third-index exclusions in the trace source do not use the old cubic, the
weight, or the carrier center.  Before assuming any sign for
\(B,C_0,R,b,V\), their outputs are
\begin{gather*}
G=g^2,\quad E=\tau^2,\quad F_0=f_0^2,\quad
z=\varepsilon gy,\quad QV^2-(Q-1)y^2=1,\\
c=\psi_A(p),\quad p\ge5\text{ odd},\quad
k=2\psi_{1+2v}(m),\quad m\ge2,\\
c>8XY^3\ge8q^4,\quad c>X,\quad K>0,\quad
b\ne0,\quad \operatorname{sgn}V=\operatorname{sgn}b,\\
Q>f^2,\quad c<f/2,\quad |V|>8Q^2,\quad
|b|>8f^3-\tfrac12.
\end{gather*}
Here \(g,\tau,f_0,y\) are positive integers and
\(\varepsilon\in\{1,-1\}\).  The sequence \(\psi_A\) has initial
values 0,1 and recurrence \(\psi_A(j+2)=2A\psi_A(j+1)-\psi_A(j)\).
If the auxiliary index were 3, its two sign cases would give, respectively,
\[
f(b-4\Delta f)=c-(4\Delta+3),\qquad
f(b+4\Delta f)=c+(4\Delta+3):
\]
Both right sides lie strictly between 0 and \(2c<f\), whereas each
left side is a multiple of $f$, a contradiction.

We now establish the new bound needed for the center.  The identities give
\[
\frac{a^2}{g^2}=\Delta U^2\frac{a}{1+a}>9,
\qquad \frac{C^2}{G}<\frac{A^2}{A^2-1}<4.
\]
Thus \(g<f/6\) and \(C<2g\).  Also
\(y<\tfrac32|V|\) and \(|V|<\tfrac32f|b|\).  For either root sign,
\begin{align*}
|\varepsilon gy+CV|
&\le gy+C|V|<\tfrac72g|V|
 <\tfrac78 f^2|b|\\
&<\tfrac7{32}b^2<\tfrac14b^2.
\end{align*}
We used \(|b|>4f^2\), a consequence of the stronger bound above.

The inverse-Lucas correction is nonnegative before rank sign recovery:
\(vk(k^2-4)/6\ge0\), because \(k\ge2\).  Consequently
\[
e_*\le k-R-1<b,\qquad \rho<\frac{|b|}{q^6}<\tfrac14b^2.
\]
If \(B\le-1\), then \(Wt<-L\).  Solving
\eqref{eq:compact-carrier} for positive \(n\) gives
\[
n<\tfrac14L-L+\tfrac14L<0,
\]
a contradiction.  Therefore \(B\ge0\), and then \(C_0>0\).
This treats all integer ranks \(R\), including negative and zero ranks.
The missing adverse-cubic estimate identified in the trace source's open
\(b^2\) question is unnecessary for the inverse correction.

\subsection{Soundness and the prime 3 ordering}

After recovering \(C_0>0\), the calibrated mask bound gives
\(3q<R<q^4-2\), hence \(c>2R\) and \(b,V>0\).  Since
\[
f-1=2\Delta i^2c^4E,\qquad V\equiv-c\pmod f,
\qquad V\equiv-R\pmod c,
\]
the same normalized-rank theorem gives \(p=R\).  In particular \(R\)
is positive odd, \(q\) is even, and \(J\) is odd.  The new center
\(S\) is an integer polynomial, so integrality of \(z\) implies
integrality of \(z+S\).  The remaining carrier argument is therefore
the one in the inverse-Lucas source; the weight enters only through
\(W\equiv1\pmod h\).

For clarity, its essential arithmetic is the recurrence expansion
\[
k=2m+8vC_m+32v^2Z_m,
\qquad C_m=\binom{m+1}{3},\qquad Z_m\in\mathbb Z_{\ge0}.
\]
Writing \(d=2m-R-1\) and \(A_m=C_m+4vZ_m\), one obtains
\[
3(e_*-d)=4v^2\bigl[24Z_m-A_m\{k^2+2mk+(2m)^2-4\}\bigr].
\]
The integer $k$ is even, so $k(k^2-4)/6$ and hence $e_*$ are integers.
Carrier integrality gives \(q^6\mid e_*-q^4s\), hence
\(q^4\mid e_*\).  Prime-by-prime valuations of the last identity give
\(q^4\mid e_*-d\), including the possible loss of one factor 3.
The bound \(1\le m<R<q^4\) forces \(d=0\).

One must not yet conclude \(q^6\mid e_*\).  First the whole-product
projection from \(H\mid6hK\), with its retained strict size bound,
gives \(X=2^R\).  Since \(q\mid X\), \(q\) is a power of 2.  Only
then can the factor 3 be canceled and \(q^6\mid e_*\) recovered.
The remaining congruences give \(q^2\mid s\), \(h\mid B\), and
integral positive \(\kappa\).  Thus
\[
X=2^R,\quad q^3\mid Y,\quad t=B/h\in\mathbb Z_{\ge0},\quad
(K_0+w)C_0-F=t(q-1).
\]
The inherited whole-tuple native/refutation calibration now yields a
finite syntactic ZFC refutation.  This step retains the entire calibrated
tuple, not merely its elementary coefficient inequalities.

\subsection{A cofinal positive lift of the centered equation}

Conversely a finite refutation supplies the canonical outer witnesses in
the two pinned source theorems.  In particular \(q\) is a power of 2,
\(R\equiv3\pmod4\), \(X=2^R\), \(q^3\mid Y\), \(C_0>0\),
and \(t\ge0\), while \(\rho,\sigma,\kappa\) are integers.
Their exact Pell-ratio interval supplies positive \(\eta,\zeta\).
Choose the strong seed with \(c^2\tau\mid I_0\), so
\(i=I_0/(c^2\tau)>0\), \(g=\tau f_0\), and \(f=2f_0^2-1\).

The retained adjacent-Pell completeness theorem supplies positive
\(V_*,y\) such that \(V_*+c+Rf^2\) is divisible by \(cf^2\).
The unsquared-trace coordinate
\[
T=\frac{V_*+c+Rf}{cf}
\]
is a positive integer, because \(c\mid f-1\).  For an uncentered
control take \(n=gy+Wt+\rho-\kappa-\sigma\).  The inherited bounds
\(|\rho|<g/3\) and \(\kappa+\sigma<g/2\) prove positivity.

For the centered equation the negative root is useful.  Put
\(u_0=U^2E\ge2\).  The exact identity
\[
(Q-1)C^2-QG=E\bigl[Q(u_0-1)-A^2u_0\bigr]
\]
is positive: \(Q>4\Delta^2u_0\) and \(4\Delta^2>A^2\).
Therefore \(C>g\sqrt{Q/(Q-1)}\).  If
\(D_0=(Q-1)C^2-QG\), then \(D_0\) is a positive integer and
\[
CV_*-gy>\frac{D_0V_*}{2C(Q-1)}\ge\frac{V_*}{2C(Q-1)}.
\]
There are cofinally many permissible \(V_*\).  The matrix
\[
\begin{pmatrix}2Q-1&2(Q-1)\\2Q&2Q-1\end{pmatrix}
\]
preserves the adjacent Pell norm, has determinant 1, and strictly
increases positive pairs.  A power of it is the identity modulo
\(cf^2\); iterating that power preserves the required congruence and
makes \(V_*\) unbounded.  No additional polynomial variable is used.

Choose \(V_*>2C(Q-1)(\kappa+\sigma+|\rho|)\) and set
\[
n=CV_*-gy+(1+hb^2)t+\rho-\kappa-\sigma.
\]
Then \(n\) is a strictly positive integer, the centered carrier is
\(z=-gy\), and \eqref{eq:compact-centered} vanishes.  This proves
positive zero-existence equivalence; it does not mistake an arbitrary
integer shear for an automorphism of the positive domain.

\subsection{Exact degree and independent finite verification}

The cancellation is an identity over a free polynomial ring:
\begin{align*}
P={}&(Q-1)Z_0^2-2(Q-1)MCVZ_0\\
&+M^2V^2\bigl[QE(U^2E-1)-C^2\bigr]+M^2G.
\end{align*}
The degrees of \(Q,M,V,U,E,Z_0,C,G\) are respectively
\(80,11,40,7,14,21,25,50\).  Its unique degree 224 term is
\(QM^2V^2U^2E^2\); all other summands have strictly lower degree.
With \(q_0=\beta J\), \(k_0=\eta+\zeta\), and
\(B_0=k_0sT-q_0^2(q_0-F)\), the complete highest homogeneous part is
\[
[P]_{224}=9\,2^{10}q_0^{96}w^{24}s^{56}i^{10}k_0^{32}B_0^2.
\]
The coefficient of
\(J^{96}w^{24}s^{58}i^{10}\eta^{34}T^2\) is
\(9\,2^{10}\beta^{96}>0\), proving the exact degree on the authentic
coefficient slice.  The uncentered controls have complete top
\[
[P]_{232}=-9\,2^{10}q_0^{100}w^{26}s^{58}i^{10}k_0^{32}B_0^2.
\]

The checker imports no producer.  It verifies every actual DAG gate by
independently specified polynomial cuts, all live inputs and coefficient
ports, the exact cancellation identity, and the complete leading forms.
Both alternative operators at every gate are deliberately introduced and
rejected:180 controls for the centered circuit and170 for each control.
Normal and optimized Python runs pass.  This authenticates the finite
artifacts; the unbounded mathematical arguments above remain essential.


\section{Verification, reproducibility, and proof boundaries}
\label{sec:verification}

The deliverable has four distinct layers of evidence.
\begin{enumerate}
\item The unbounded pooling theorem, its zero-base bootstrap, the compact
transport estimate, and the cofinal positive lift are mathematical proofs.
The independent internal reviews separately check these changed arguments.
\item The frozen parent, full request table, and literal pool expansions
are authenticated with exact integer arithmetic. The checker reconstructs
all seventeen request energies and the retained Pell and multiplication
rows, checks coordinate privacy, and verifies every output coefficient.
\item The projection checker imports no producer or optimizer. It
independently checks unit pivots, acyclicity, every local circuit block,
the complete source-row partition, factored pool formulas, the final sum
of squares, exact degrees, coefficient bounds, and input dependence.
Corrupted artifacts are deliberately rejected.
\item The compact checker verifies all polynomial circuit gates through
independent cuts, the root-center cancellation, the complete leading
forms, and every fixed coefficient payload. Its constructor reproduces
the four coefficient integers from the printed finite seed and recipe.
\end{enumerate}

The retained parent theorem and the named C222/ZF/ZFC logical calibration
are inherited mathematical dependencies. This package does not contain a
new transitive Lean or Rocq build, does not emit an encoded object-language
ZFC proof, and does not prove $\Con(\mathrm{ZFC})$. Independent AI-assisted
reviews are not external human peer review. The ordinary equivalence
proof and the finite algebraic verification should be evaluated at their
respective scopes.

All core replay programs use Python's standard library. From the package
root run
\begin{verbatim}
python3 -B verify.py
python3 -B -O verify.py
python3 -B code/build_projections.py --check
\end{verbatim}
The first command verifies the saved deliverable without performing a
root search. The evaluator accepts an explicit tuple and computes its
polynomial value or its residuals; it is not a consistency decision
procedure. The chosen pivot lists are frozen certificates; rebuilding and
verifying them does not require an optimizer or repeat the exploratory search.

The article can be rebuilt with two runs of pdfLaTeX. The exact
mathematical definitions are the source coefficient files, pivot
certificates, factored circuits, and printed appendices; PDF bytes can
vary with the TeX installation. Source pins, attribution, licenses,
and a release checksum manifest accompany the package.

\section{Further research questions}

\paragraph{Choose the best partition of the bases.}
The all-five, staged, and paired formulas illustrate a genuine degree
versus coordinate tradeoff. The best grouping can depend on which semantic
coordinates a later monic elimination replaces. An optimizer that jointly
chooses groups and polynomial definitions could improve on choosing the
grouping first. Its output should remain a small exact certificate.

\paragraph{Compress the signed Pell frame itself.}
The present theorem shares power-congruence certificates while retaining
the parent's canonical-index frame. The next substantial saving would
compress those index and orientation witnesses without assuming a desired
sign or rank prematurely. A new shared identity should be tested first
against negative indices, unit bases, and temporary zero bases.

\paragraph{Exploit stronger height slack.}
The proof only uses $W^{6b+2}\le W^W$ with $b\ge16$, leaving a large
margin. Additional bounded arithmetic tests might fit into the same
rational energy. A proposal must preserve the common sign and justify
how every modulus remains nonzero, rather than merely multiplying more
congruences together.

\paragraph{Lower the degree of the smallest direct forms.}
The highest degrees in the 87-, 88-, and 89-variable forms come from long
product and squaring chains, not solely from the five-base numerator.
Introducing a few strategically selected intermediates, or changing the
factorization of the Horn-flow certificate, may lower degree at a better
coordinate cost. None of the present pivot lists is claimed optimal.

\paragraph{Shorten the centered compact evaluator.}
The new degree-224 carrier uses five additional operations compared with
the 85-operation controls. A different schedule for the root center, or
a different polynomial approximation to $\sqrt G$, could improve that
tradeoff. Positive completeness must be proved for any new center;
an arbitrary integer shear need not preserve positive witnesses.

\paragraph{Formalize the new interface.}
A useful focused formalization target is Theorem~\ref{thm:all-bases}
with an explicit list of bases and requests, followed by its literal
17-request instance. This would make future groupings reuse a single
signed, zero-base-aware theorem. The complete consistency statement would
then compose with the repository's existing formal calibration interfaces.

\begin{thebibliography}{99}
\raggedright
\bibitem{parent158} ProveIt, \emph{158-pooled-powers}, especially
\src{equation/equation.json}, \src{proofs/POOLED_POWERS.md},
\src{proofs/STRICT_HORIZON_AND_BETA.md}, and the retained bootstrap proof.
Snapshot \texttt{0c48d332df6b0e50d4e5b6251c531898f5682c72}.
\url{https://github.com/VladimirReshetnikov/ProveIt/tree/0c48d332df6b0e50d4e5b6251c531898f5682c72/Computability/DiophantineConsistency/con-zfc-diophantine-forms/158-pooled-powers}.
\bibitem{pool168} ProveIt, \emph{168-reciprocal-pooling},
\src{proofs/RECIPROCAL_POOLING.md}, same snapshot.
\url{https://github.com/VladimirReshetnikov/ProveIt/tree/0c48d332df6b0e50d4e5b6251c531898f5682c72/Computability/DiophantineConsistency/con-zfc-diophantine-forms/168-reciprocal-pooling}.
\bibitem{projection152} ProveIt, \emph{152-lucas-corrections},
\src{higher_degree/section.tex} and \src{higher_degree/RESULTS.json}, same snapshot.
\url{https://github.com/VladimirReshetnikov/ProveIt/tree/0c48d332df6b0e50d4e5b6251c531898f5682c72/Computability/DiophantineConsistency/con-zfc-diophantine-forms/152-lucas-corrections}.
\bibitem{trace163} ProveIt, \emph{163-trace-radical},
\src{compact/PROOF.md} and its complete coefficient recipe, same snapshot.
\url{https://github.com/VladimirReshetnikov/ProveIt/tree/0c48d332df6b0e50d4e5b6251c531898f5682c72/Computability/DiophantineConsistency/con-zfc-diophantine-forms/163-trace-radical}.
\bibitem{inverse164} ProveIt, \emph{164-unsquared-factors-linear-selectors},
\src{positive/PROOF.md} and the calibrated parent circuit, same snapshot.
\url{https://github.com/VladimirReshetnikov/ProveIt/tree/0c48d332df6b0e50d4e5b6251c531898f5682c72/Computability/DiophantineConsistency/con-zfc-diophantine-forms/164-unsquared-factors-linear-selectors}.
\bibitem{calibration167} ProveIt, \emph{167-integral-projection},
\src{proofs/INHERITED_PROOF_AUDIT.md}, same snapshot.
\url{https://github.com/VladimirReshetnikov/ProveIt/tree/0c48d332df6b0e50d4e5b6251c531898f5682c72/Computability/DiophantineConsistency/con-zfc-diophantine-forms/167-integral-projection}.
\bibitem{matiyasevich} Yu. Matiyasevich, \emph{Diophantine representation
of enumerable predicates}, Izvestiya Akademii Nauk SSSR, Seriya
Matematicheskaya 35 (1971), 3--30.
\url{https://geodesic.mathdoc.fr/item/IM2_1971_5_1_a0/}.
\bibitem{matbook} Yu. Matiyasevich, \emph{Hilbert's Tenth Problem},
MIT Press, 1993. Author's online commentary to Chapter 5:
\url{https://logic.pdmi.ras.ru/~yumat/H10Pbook/commch_5.htm}.
\bibitem{sun} Zhi-Wei Sun, \emph{Ten unknowns for Hilbert's tenth problem
over the integers}, arXiv:2609.01594v2, 2 September 2026.
\url{https://arxiv.org/abs/2609.01594v2}.
\end{thebibliography}

\appendix
\section{Every residual of the authenticated parent}
\label{app:parent}
All variables $x_0,\ldots,x_{224}$ in this appendix range over unrestricted
integers. The polynomials below are literally the rows of the frozen
parent; no omitted scalar parameter occurs. Each expression defines
$R_r$, and the parent equation is $\sum_{r=0}^{161}R_r^2=0$.
The main text gives the exact coordinate and row deletion lists, the
complete request table, and the new pooling rows. These together specify
the three raw systems used by the next appendix.
\begingroup\small
%BEGIN_RESIDUAL 0
\[
\begin{aligned}
R_{0}={}&-x_{2}^{2}+x_{3}
\end{aligned}
\]
%END_RESIDUAL 0
%BEGIN_RESIDUAL 1
\[
\begin{aligned}
R_{1}={}&x_{0}-x_{0}x_{5}-x_{5}+x_{8}
\end{aligned}
\]
%END_RESIDUAL 1
%BEGIN_RESIDUAL 2
\[
\begin{aligned}
R_{2}={}&-1+x_{1}-x_{2}-x_{5}x_{6}
\end{aligned}
\]
%END_RESIDUAL 2
%BEGIN_RESIDUAL 3
\[
\begin{aligned}
R_{3}={}&-1+x_{1}-x_{7}x_{8}
\end{aligned}
\]
%END_RESIDUAL 3
%BEGIN_RESIDUAL 4
\[
\begin{aligned}
R_{4}={}&1+x_{3}-x_{4}x_{19}+x_{17}
\end{aligned}
\]
%END_RESIDUAL 4
%BEGIN_RESIDUAL 5
\[
\begin{aligned}
R_{5}={}&-1+2x_{9}-x_{9}^{2}-2x_{9}x_{10}+2x_{10}-x_{10}^{2}+x_{20}
\end{aligned}
\]
%END_RESIDUAL 5
%BEGIN_RESIDUAL 6
\[
\begin{aligned}
R_{6}={}&-5-30x_{9}-36x_{9}^{2}-96x_{9}x_{10}-24x_{9}x_{20}-40x_{10}-64x_{10}^{2}-32x_{10}x_{20}\\
&-10x_{20}-4x_{20}^{2}+x_{21}
\end{aligned}
\]
%END_RESIDUAL 6
%BEGIN_RESIDUAL 7
\[
\begin{aligned}
R_{7}={}&-2+2x_{9}-x_{9}^{2}-4x_{9}x_{21}+4x_{21}-4x_{21}^{2}+x_{22}
\end{aligned}
\]
%END_RESIDUAL 7
%BEGIN_RESIDUAL 8
\[
\begin{aligned}
R_{8}={}&1+6x_{9}-36x_{9}^{2}-192x_{9}x_{21}-24x_{9}x_{22}+16x_{21}-256x_{21}^{2}-64x_{21}x_{22}\\
&+2x_{22}-4x_{22}^{2}+x_{23}
\end{aligned}
\]
%END_RESIDUAL 8
%BEGIN_RESIDUAL 9
\[
\begin{aligned}
R_{9}={}&-2-x_{10}^{2}-2x_{10}x_{11}-x_{11}^{2}+x_{24}
\end{aligned}
\]
%END_RESIDUAL 9
%BEGIN_RESIDUAL 10
\[
\begin{aligned}
R_{10}={}&-9+7x_{9}-x_{9}^{2}-4x_{9}x_{10}-4x_{9}x_{24}+12x_{10}-4x_{10}^{2}-8x_{10}x_{24}\\
&+12x_{24}-4x_{24}^{2}+x_{25}
\end{aligned}
\]
%END_RESIDUAL 10
%BEGIN_RESIDUAL 11
\[
\begin{aligned}
R_{11}={}&-1+2x_{9}-x_{9}^{2}-2x_{9}x_{11}+2x_{11}-x_{11}^{2}+x_{26}
\end{aligned}
\]
%END_RESIDUAL 11
%BEGIN_RESIDUAL 12
\[
\begin{aligned}
R_{12}={}&-6x_{9}-144x_{9}^{2}-192x_{9}x_{10}-192x_{9}x_{11}-48x_{9}x_{20}-48x_{9}x_{26}-8x_{10}-64x_{10}^{2}\\
&-128x_{10}x_{11}-32x_{10}x_{20}-32x_{10}x_{26}-64x_{11}^{2}-32x_{11}x_{20}-32x_{11}x_{26}-2x_{20}-4x_{20}^{2}\\
&-8x_{20}x_{26}-4x_{26}^{2}+x_{27}
\end{aligned}
\]
%END_RESIDUAL 12
%BEGIN_RESIDUAL 13
\[
\begin{aligned}
R_{13}={}&-4x_{25}^{2}-8x_{25}x_{27}-4x_{27}^{2}+x_{28}
\end{aligned}
\]
%END_RESIDUAL 13
%BEGIN_RESIDUAL 14
\[
\begin{aligned}
R_{14}={}&-6+16x_{11}-16x_{11}^{2}-16x_{11}x_{24}-16x_{11}x_{26}+8x_{24}-4x_{24}^{2}-8x_{24}x_{26}\\
&+8x_{26}-4x_{26}^{2}+x_{29}
\end{aligned}
\]
%END_RESIDUAL 14
%BEGIN_RESIDUAL 15
\[
\begin{aligned}
R_{15}={}&1+6x_{9}-36x_{9}^{2}-96x_{9}x_{10}-48x_{9}x_{11}-24x_{9}x_{20}-48x_{9}x_{26}-24x_{9}x_{29}\\
&+8x_{10}-64x_{10}^{2}-64x_{10}x_{11}-32x_{10}x_{20}-64x_{10}x_{26}-32x_{10}x_{29}+8x_{11}-16x_{11}^{2}\\
&-16x_{11}x_{20}-32x_{11}x_{26}-16x_{11}x_{29}+2x_{20}-4x_{20}^{2}-16x_{20}x_{26}-8x_{20}x_{29}+8x_{26}\\
&-16x_{26}^{2}-16x_{26}x_{29}+4x_{29}-4x_{29}^{2}+x_{30}
\end{aligned}
\]
%END_RESIDUAL 15
%BEGIN_RESIDUAL 16
\[
\begin{aligned}
R_{16}={}&-10-5x_{9}-x_{9}^{2}-12x_{9}x_{10}-16x_{9}x_{11}-4x_{9}x_{24}-36x_{10}-36x_{10}^{2}\\
&-96x_{10}x_{11}-24x_{10}x_{24}-48x_{11}-64x_{11}^{2}-32x_{11}x_{24}-12x_{24}-4x_{24}^{2}+x_{31}
\end{aligned}
\]
%END_RESIDUAL 16
%BEGIN_RESIDUAL 17
\[
\begin{aligned}
R_{17}={}&-2-24x_{9}-36x_{9}^{2}-12x_{9}x_{10}-96x_{9}x_{11}-24x_{9}x_{26}-3x_{10}-x_{10}^{2}\\
&-16x_{10}x_{11}-4x_{10}x_{26}-32x_{11}-64x_{11}^{2}-32x_{11}x_{26}-8x_{26}-4x_{26}^{2}+x_{32}
\end{aligned}
\]
%END_RESIDUAL 17
%BEGIN_RESIDUAL 18
\[
\begin{aligned}
R_{18}={}&2-4x_{31}^{2}-8x_{31}x_{32}-4x_{32}^{2}+x_{33}
\end{aligned}
\]
%END_RESIDUAL 18
%BEGIN_RESIDUAL 19
\[
\begin{aligned}
R_{19}={}&-12+16x_{9}-4x_{9}^{2}-8x_{9}x_{11}-8x_{9}x_{20}-8x_{9}x_{24}+16x_{11}-4x_{11}^{2}\\
&-8x_{11}x_{20}-8x_{11}x_{24}+16x_{20}-4x_{20}^{2}-8x_{20}x_{24}+16x_{24}-4x_{24}^{2}+x_{34}
\end{aligned}
\]
%END_RESIDUAL 19
%BEGIN_RESIDUAL 20
\[
\begin{aligned}
R_{20}={}&-3-24x_{9}-64x_{9}^{2}-160x_{9}x_{11}-64x_{9}x_{24}-32x_{9}x_{26}-32x_{9}x_{34}-26x_{11}\\
&-100x_{11}^{2}-80x_{11}x_{24}-40x_{11}x_{26}-40x_{11}x_{34}-8x_{24}-16x_{24}^{2}-16x_{24}x_{26}-16x_{24}x_{34}\\
&-6x_{26}-4x_{26}^{2}-8x_{26}x_{34}-4x_{34}-4x_{34}^{2}+x_{35}
\end{aligned}
\]
%END_RESIDUAL 20
%BEGIN_RESIDUAL 21
\[
\begin{aligned}
R_{21}={}&-1-18x_{9}-36x_{9}^{2}-96x_{9}x_{10}-24x_{9}x_{20}-24x_{10}-64x_{10}^{2}-32x_{10}x_{20}\\
&-6x_{20}-4x_{20}^{2}+x_{36}
\end{aligned}
\]
%END_RESIDUAL 21
%BEGIN_RESIDUAL 22
\[
\begin{aligned}
R_{22}={}&-2+2x_{9}-x_{9}^{2}-4x_{9}x_{36}+4x_{36}-4x_{36}^{2}+x_{37}
\end{aligned}
\]
%END_RESIDUAL 22
%BEGIN_RESIDUAL 23
\[
\begin{aligned}
R_{23}={}&1+6x_{9}-36x_{9}^{2}-192x_{9}x_{36}-24x_{9}x_{37}+16x_{36}-256x_{36}^{2}-64x_{36}x_{37}\\
&+2x_{37}-4x_{37}^{2}+x_{38}
\end{aligned}
\]
%END_RESIDUAL 23
%BEGIN_RESIDUAL 24
\[
\begin{aligned}
R_{24}={}&-x_{10}^{2}-4x_{10}x_{38}-4x_{38}^{2}+x_{39}
\end{aligned}
\]
%END_RESIDUAL 24
%BEGIN_RESIDUAL 25
\[
\begin{aligned}
R_{25}={}&-6-30x_{9}-64x_{9}^{2}-128x_{9}x_{10}-64x_{9}x_{20}-32x_{10}-64x_{10}^{2}-64x_{10}x_{20}\\
&-14x_{20}-16x_{20}^{2}+x_{40}
\end{aligned}
\]
%END_RESIDUAL 25
%BEGIN_RESIDUAL 26
\[
\begin{aligned}
R_{26}={}&-4+8x_{10}-64x_{10}^{2}-96x_{10}x_{11}-32x_{10}x_{24}-32x_{10}x_{40}+6x_{11}-36x_{11}^{2}\\
&-24x_{11}x_{24}-24x_{11}x_{40}+2x_{24}-4x_{24}^{2}-8x_{24}x_{40}+4x_{40}-4x_{40}^{2}+x_{41}
\end{aligned}
\]
%END_RESIDUAL 26
%BEGIN_RESIDUAL 27
\[
\begin{aligned}
R_{27}={}&8x_{9}-64x_{9}^{2}-96x_{9}x_{11}-32x_{9}x_{26}-32x_{9}x_{41}+6x_{11}-36x_{11}^{2}-24x_{11}x_{26}\\
&-24x_{11}x_{41}+2x_{26}-4x_{26}^{2}-8x_{26}x_{41}+4x_{41}-4x_{41}^{2}+x_{42}
\end{aligned}
\]
%END_RESIDUAL 27
%BEGIN_RESIDUAL 28
\[
\begin{aligned}
R_{28}={}&-x_{11}^{2}-2x_{11}x_{12}-x_{12}^{2}+x_{43}
\end{aligned}
\]
%END_RESIDUAL 28
%BEGIN_RESIDUAL 29
\[
\begin{aligned}
R_{29}={}&2x_{9}-x_{9}^{2}-4x_{9}x_{11}-4x_{9}x_{43}+4x_{11}-4x_{11}^{2}-8x_{11}x_{43}+4x_{43}\\
&-4x_{43}^{2}+x_{44}
\end{aligned}
\]
%END_RESIDUAL 29
%BEGIN_RESIDUAL 30
\[
\begin{aligned}
R_{30}={}&-35+24x_{9}-4x_{9}^{2}-8x_{9}x_{10}-8x_{9}x_{24}-8x_{9}x_{26}+24x_{10}-4x_{10}^{2}\\
&-8x_{10}x_{24}-8x_{10}x_{26}+24x_{24}-4x_{24}^{2}-8x_{24}x_{26}+24x_{26}-4x_{26}^{2}+x_{45}
\end{aligned}
\]
%END_RESIDUAL 30
%BEGIN_RESIDUAL 31
\[
\begin{aligned}
R_{31}={}&-16x_{9}^{2}-32x_{9}x_{11}-32x_{9}x_{26}-32x_{9}x_{43}-16x_{9}x_{44}-16x_{9}x_{45}-4x_{11}-16x_{11}^{2}\\
&-32x_{11}x_{26}-32x_{11}x_{43}-16x_{11}x_{44}-16x_{11}x_{45}-16x_{26}^{2}-32x_{26}x_{43}-16x_{26}x_{44}\\
&-16x_{26}x_{45}-4x_{43}-16x_{43}^{2}-16x_{43}x_{44}-16x_{43}x_{45}-2x_{44}-4x_{44}^{2}-8x_{44}x_{45}\\
&-4x_{45}^{2}+x_{46}
\end{aligned}
\]
%END_RESIDUAL 31
%BEGIN_RESIDUAL 32
\[
\begin{aligned}
R_{32}={}&x_{11}-25x_{11}^{2}-40x_{11}x_{26}-20x_{11}x_{29}-16x_{26}^{2}-16x_{26}x_{29}-4x_{29}^{2}+x_{47}
\end{aligned}
\]
%END_RESIDUAL 32
%BEGIN_RESIDUAL 33
\[
\begin{aligned}
R_{33}={}&-1+8x_{10}-4x_{10}^{2}-8x_{10}x_{24}+8x_{24}-4x_{24}^{2}+x_{48}
\end{aligned}
\]
%END_RESIDUAL 33
%BEGIN_RESIDUAL 34
\[
\begin{aligned}
R_{34}={}&-1-16x_{10}^{2}-8x_{10}x_{11}-32x_{10}x_{24}-16x_{10}x_{48}-x_{11}^{2}-8x_{11}x_{24}-4x_{11}x_{48}\\
&-16x_{24}^{2}-16x_{24}x_{48}-4x_{48}^{2}+x_{49}
\end{aligned}
\]
%END_RESIDUAL 34
%BEGIN_RESIDUAL 35
\[
\begin{aligned}
R_{35}={}&1+x_{10}-x_{10}^{2}-4x_{10}x_{11}-4x_{10}x_{49}-4x_{11}^{2}-8x_{11}x_{49}-4x_{49}^{2}\\
&+x_{50}
\end{aligned}
\]
%END_RESIDUAL 35
%BEGIN_RESIDUAL 36
\[
\begin{aligned}
R_{36}={}&-1-4x_{50}-4x_{50}^{2}+x_{51}
\end{aligned}
\]
%END_RESIDUAL 36
%BEGIN_RESIDUAL 37
\[
\begin{aligned}
R_{37}={}&-1-x_{11}^{2}-8x_{11}x_{50}-4x_{11}x_{51}-16x_{50}^{2}-16x_{50}x_{51}-4x_{51}^{2}+x_{52}
\end{aligned}
\]
%END_RESIDUAL 37
%BEGIN_RESIDUAL 38
\[
\begin{aligned}
R_{38}={}&2+x_{10}-x_{10}^{2}-4x_{10}x_{11}-4x_{10}x_{52}-4x_{11}^{2}-8x_{11}x_{52}-4x_{52}^{2}\\
&+x_{53}
\end{aligned}
\]
%END_RESIDUAL 38
%BEGIN_RESIDUAL 39
\[
\begin{aligned}
R_{39}={}&-4x_{47}^{2}-8x_{47}x_{53}-4x_{53}^{2}+x_{54}
\end{aligned}
\]
%END_RESIDUAL 39
%BEGIN_RESIDUAL 40
\[
\begin{aligned}
R_{40}={}&-9+7x_{9}-x_{9}^{2}-8x_{9}x_{47}-4x_{9}x_{54}+24x_{47}-16x_{47}^{2}-16x_{47}x_{54}\\
&+12x_{54}-4x_{54}^{2}+x_{55}
\end{aligned}
\]
%END_RESIDUAL 40
%BEGIN_RESIDUAL 41
\[
\begin{aligned}
R_{41}={}&-1-4x_{55}-4x_{55}^{2}+x_{56}
\end{aligned}
\]
%END_RESIDUAL 41
%BEGIN_RESIDUAL 42
\[
\begin{aligned}
R_{42}={}&-x_{11}^{2}-8x_{11}x_{55}-4x_{11}x_{56}-16x_{55}^{2}-16x_{55}x_{56}-4x_{56}^{2}+x_{57}
\end{aligned}
\]
%END_RESIDUAL 42
%BEGIN_RESIDUAL 43
\[
\begin{aligned}
R_{43}={}&-x_{9}-x_{9}^{2}-4x_{9}x_{11}-4x_{9}x_{57}-4x_{11}-4x_{11}^{2}-8x_{11}x_{57}-4x_{57}\\
&-4x_{57}^{2}+x_{58}
\end{aligned}
\]
%END_RESIDUAL 43
%BEGIN_RESIDUAL 44
\[
\begin{aligned}
R_{44}={}&-10-12x_{58}-4x_{58}^{2}+x_{59}
\end{aligned}
\]
%END_RESIDUAL 44
%BEGIN_RESIDUAL 45
\[
\begin{aligned}
R_{45}={}&1+4x_{10}-4x_{10}^{2}-8x_{10}x_{24}-16x_{10}x_{58}-8x_{10}x_{59}+4x_{24}-4x_{24}^{2}\\
&-16x_{24}x_{58}-8x_{24}x_{59}+8x_{58}-16x_{58}^{2}-16x_{58}x_{59}+4x_{59}-4x_{59}^{2}+x_{60}
\end{aligned}
\]
%END_RESIDUAL 45
%BEGIN_RESIDUAL 46
\[
\begin{aligned}
R_{46}={}&x_{10}-25x_{10}^{2}-40x_{10}x_{24}-20x_{10}x_{60}-16x_{24}^{2}-16x_{24}x_{60}-4x_{60}^{2}+x_{61}
\end{aligned}
\]
%END_RESIDUAL 46
%BEGIN_RESIDUAL 47
\[
\begin{aligned}
R_{47}={}&-10-12x_{53}-4x_{53}^{2}+x_{62}
\end{aligned}
\]
%END_RESIDUAL 47
%BEGIN_RESIDUAL 48
\[
\begin{aligned}
R_{48}={}&1-16x_{53}^{2}-16x_{53}x_{61}-16x_{53}x_{62}-4x_{61}^{2}-8x_{61}x_{62}-4x_{62}^{2}+x_{63}
\end{aligned}
\]
%END_RESIDUAL 48
%BEGIN_RESIDUAL 49
\[
\begin{aligned}
R_{49}={}&x_{11}-x_{11}^{2}-8x_{11}x_{61}-4x_{11}x_{63}-16x_{61}^{2}-16x_{61}x_{63}-4x_{63}^{2}+x_{64}
\end{aligned}
\]
%END_RESIDUAL 49
%BEGIN_RESIDUAL 50
\[
\begin{aligned}
R_{50}={}&-4x_{9}-4x_{9}^{2}-4x_{9}x_{11}-8x_{9}x_{20}-x_{11}-x_{11}^{2}-4x_{11}x_{20}-4x_{20}\\
&-4x_{20}^{2}+x_{65}
\end{aligned}
\]
%END_RESIDUAL 50
%BEGIN_RESIDUAL 51
\[
\begin{aligned}
R_{51}={}&1+4x_{9}-4x_{9}^{2}-8x_{9}x_{20}-8x_{9}x_{65}+4x_{20}-4x_{20}^{2}-8x_{20}x_{65}\\
&+4x_{65}-4x_{65}^{2}+x_{66}
\end{aligned}
\]
%END_RESIDUAL 51
%BEGIN_RESIDUAL 52
\[
\begin{aligned}
R_{52}={}&-2-4x_{9}-16x_{9}^{2}-32x_{9}x_{20}-16x_{9}x_{64}-16x_{9}x_{66}-4x_{20}-16x_{20}^{2}\\
&-16x_{20}x_{64}-16x_{20}x_{66}-4x_{64}^{2}-8x_{64}x_{66}-2x_{66}-4x_{66}^{2}+x_{67}
\end{aligned}
\]
%END_RESIDUAL 52
%BEGIN_RESIDUAL 53
\[
\begin{aligned}
R_{53}={}&-2-4x_{67}-4x_{67}^{2}+x_{68}
\end{aligned}
\]
%END_RESIDUAL 53
%BEGIN_RESIDUAL 54
\[
\begin{aligned}
R_{54}={}&-x_{2}x_{14}+x_{69}
\end{aligned}
\]
%END_RESIDUAL 54
%BEGIN_RESIDUAL 55
\[
\begin{aligned}
R_{55}={}&x_{13}-x_{14}-2x_{23}x_{69}-x_{70}
\end{aligned}
\]
%END_RESIDUAL 55
%BEGIN_RESIDUAL 56
\[
\begin{aligned}
R_{56}={}&x_{13}-x_{14}-4x_{25}x_{69}-2x_{28}x_{69}-x_{71}
\end{aligned}
\]
%END_RESIDUAL 56
%BEGIN_RESIDUAL 57
\[
\begin{aligned}
R_{57}={}&x_{13}-x_{14}-2x_{30}x_{69}-6x_{69}-x_{72}
\end{aligned}
\]
%END_RESIDUAL 57
%BEGIN_RESIDUAL 58
\[
\begin{aligned}
R_{58}={}&x_{13}-x_{14}-4x_{31}x_{69}-2x_{33}x_{69}-x_{73}
\end{aligned}
\]
%END_RESIDUAL 58
%BEGIN_RESIDUAL 59
\[
\begin{aligned}
R_{59}={}&x_{13}-x_{14}-2x_{35}x_{69}+6x_{69}-x_{74}
\end{aligned}
\]
%END_RESIDUAL 59
%BEGIN_RESIDUAL 60
\[
\begin{aligned}
R_{60}={}&x_{13}-x_{14}-4x_{38}x_{69}-2x_{39}x_{69}-x_{75}
\end{aligned}
\]
%END_RESIDUAL 60
%BEGIN_RESIDUAL 61
\[
\begin{aligned}
R_{61}={}&x_{13}-x_{14}-2x_{42}x_{69}-x_{76}
\end{aligned}
\]
%END_RESIDUAL 61
%BEGIN_RESIDUAL 62
\[
\begin{aligned}
R_{62}={}&x_{13}-x_{14}-2x_{46}x_{69}-4x_{69}-x_{77}
\end{aligned}
\]
%END_RESIDUAL 62
%BEGIN_RESIDUAL 63
\[
\begin{aligned}
R_{63}={}&-x_{9}x_{10}+x_{9}x_{11}+x_{10}+x_{10}x_{11}-x_{11}-x_{11}^{2}+x_{78}
\end{aligned}
\]
%END_RESIDUAL 63
%BEGIN_RESIDUAL 64
\[
\begin{aligned}
R_{64}={}&1+2x_{12}+x_{12}^{2}-x_{78}^{2}+x_{79}
\end{aligned}
\]
%END_RESIDUAL 64
%BEGIN_RESIDUAL 65
\[
\begin{aligned}
R_{65}={}&-x_{13}^{2}+2x_{13}x_{14}+2x_{13}x_{81}-x_{14}^{2}-2x_{14}x_{81}-x_{79}^{2}+x_{80}-x_{81}^{2}
\end{aligned}
\]
%END_RESIDUAL 65
%BEGIN_RESIDUAL 66
\[
\begin{aligned}
R_{66}={}&-4x_{67}x_{69}-2x_{68}x_{69}+2x_{69}+x_{81}
\end{aligned}
\]
%END_RESIDUAL 66
%BEGIN_RESIDUAL 67
\[
\begin{aligned}
R_{67}={}&-x_{10}x_{69}-x_{14}+x_{82}x_{83}+x_{84}
\end{aligned}
\]
%END_RESIDUAL 67
%BEGIN_RESIDUAL 68
\[
\begin{aligned}
R_{68}={}&-1+2x_{2}-2x_{2}x_{9}-2x_{2}x_{20}+x_{82}
\end{aligned}
\]
%END_RESIDUAL 68
%BEGIN_RESIDUAL 69
\[
\begin{aligned}
R_{69}={}&-x_{9}x_{220}-x_{13}+x_{83}+x_{220}
\end{aligned}
\]
%END_RESIDUAL 69
%BEGIN_RESIDUAL 70
\[
\begin{aligned}
R_{70}={}&-6x_{9}x_{69}-8x_{10}x_{69}+x_{10}x_{220}+x_{13}-x_{14}-2x_{20}x_{69}+x_{85}
\end{aligned}
\]
%END_RESIDUAL 70
%BEGIN_RESIDUAL 71
\[
\begin{aligned}
R_{71}={}&-x_{13}x_{70}+x_{14}x_{70}+x_{86}
\end{aligned}
\]
%END_RESIDUAL 71
%BEGIN_RESIDUAL 72
\[
\begin{aligned}
R_{72}={}&-x_{71}x_{86}+x_{87}
\end{aligned}
\]
%END_RESIDUAL 72
%BEGIN_RESIDUAL 73
\[
\begin{aligned}
R_{73}={}&-x_{72}x_{87}+x_{88}
\end{aligned}
\]
%END_RESIDUAL 73
%BEGIN_RESIDUAL 74
\[
\begin{aligned}
R_{74}={}&-x_{73}x_{88}+x_{89}
\end{aligned}
\]
%END_RESIDUAL 74
%BEGIN_RESIDUAL 75
\[
\begin{aligned}
R_{75}={}&-x_{74}x_{89}+x_{90}
\end{aligned}
\]
%END_RESIDUAL 75
%BEGIN_RESIDUAL 76
\[
\begin{aligned}
R_{76}={}&-x_{75}x_{90}+x_{91}
\end{aligned}
\]
%END_RESIDUAL 76
%BEGIN_RESIDUAL 77
\[
\begin{aligned}
R_{77}={}&-x_{76}x_{91}+x_{92}
\end{aligned}
\]
%END_RESIDUAL 77
%BEGIN_RESIDUAL 78
\[
\begin{aligned}
R_{78}={}&-x_{77}x_{92}+x_{93}
\end{aligned}
\]
%END_RESIDUAL 78
%BEGIN_RESIDUAL 79
\[
\begin{aligned}
R_{79}={}&-x_{84}x_{93}+x_{94}
\end{aligned}
\]
%END_RESIDUAL 79
%BEGIN_RESIDUAL 80
\[
\begin{aligned}
R_{80}={}&-x_{85}x_{94}+x_{95}
\end{aligned}
\]
%END_RESIDUAL 80
%BEGIN_RESIDUAL 81
\[
\begin{aligned}
R_{81}={}&-x_{95}^{2}+x_{96}
\end{aligned}
\]
%END_RESIDUAL 81
%BEGIN_RESIDUAL 82
\[
\begin{aligned}
R_{82}={}&-x_{5}-x_{5}x_{17}+x_{17}+x_{97}
\end{aligned}
\]
%END_RESIDUAL 82
%BEGIN_RESIDUAL 83
\[
\begin{aligned}
R_{83}={}&-x_{15}x_{97}+x_{98}
\end{aligned}
\]
%END_RESIDUAL 83
%BEGIN_RESIDUAL 84
\[
\begin{aligned}
R_{84}={}&-2x_{1}^{2}+2x_{1}x_{13}+2x_{1}x_{14}+2x_{1}x_{98}+2x_{1}x_{99}-x_{13}^{2}-2x_{13}x_{98}-x_{14}^{2}\\
&-2x_{14}x_{99}+x_{18}x_{101}-x_{80}x_{96}-x_{98}^{2}-x_{99}^{2}
\end{aligned}
\]
%END_RESIDUAL 84
%BEGIN_RESIDUAL 85
\[
\begin{aligned}
R_{85}={}&-x_{5}x_{16}+x_{16}-x_{16}x_{97}+x_{99}
\end{aligned}
\]
%END_RESIDUAL 85
%BEGIN_RESIDUAL 86
\[
\begin{aligned}
R_{86}={}&-9x_{9}-9x_{10}-4x_{20}-2x_{23}-4x_{25}-2x_{28}-2x_{30}-4x_{31}\\
&-2x_{33}-2x_{35}-4x_{38}-2x_{39}-2x_{42}-2x_{46}-4x_{67}-2x_{68}\\
&+x_{100}
\end{aligned}
\]
%END_RESIDUAL 86
%BEGIN_RESIDUAL 87
\[
\begin{aligned}
R_{87}={}&-x_{102}x_{105}+x_{105}+x_{107}
\end{aligned}
\]
%END_RESIDUAL 87
%BEGIN_RESIDUAL 88
\[
\begin{aligned}
R_{88}={}&-x_{4}x_{128}+x_{103}x_{210}-x_{106}-x_{107}x_{128}
\end{aligned}
\]
%END_RESIDUAL 88
%BEGIN_RESIDUAL 89
\[
\begin{aligned}
R_{89}={}&x_{108}x_{210}-x_{109}x_{131}-x_{110}
\end{aligned}
\]
%END_RESIDUAL 89
%BEGIN_RESIDUAL 90
\[
\begin{aligned}
R_{90}={}&x_{111}x_{210}-x_{112}x_{131}-x_{113}
\end{aligned}
\]
%END_RESIDUAL 90
%BEGIN_RESIDUAL 91
\[
\begin{aligned}
R_{91}={}&x_{114}x_{210}-x_{115}x_{131}-x_{116}
\end{aligned}
\]
%END_RESIDUAL 91
%BEGIN_RESIDUAL 92
\[
\begin{aligned}
R_{92}={}&x_{117}x_{210}-x_{118}x_{131}-x_{119}
\end{aligned}
\]
%END_RESIDUAL 92
%BEGIN_RESIDUAL 93
\[
\begin{aligned}
R_{93}={}&x_{120}x_{211}-x_{122}x_{131}-x_{123}
\end{aligned}
\]
%END_RESIDUAL 93
%BEGIN_RESIDUAL 94
\[
\begin{aligned}
R_{94}={}&x_{124}x_{211}-x_{125}x_{131}-x_{126}
\end{aligned}
\]
%END_RESIDUAL 94
%BEGIN_RESIDUAL 95
\[
\begin{aligned}
R_{95}={}&-x_{18}x_{128}+x_{131}
\end{aligned}
\]
%END_RESIDUAL 95
%BEGIN_RESIDUAL 96
\[
\begin{aligned}
R_{96}={}&x_{127}-x_{129}x_{131}-x_{130}
\end{aligned}
\]
%END_RESIDUAL 96
%BEGIN_RESIDUAL 97
\[
\begin{aligned}
R_{97}={}&-x_{131}x_{133}+x_{132}-x_{134}
\end{aligned}
\]
%END_RESIDUAL 97
%BEGIN_RESIDUAL 98
\[
\begin{aligned}
R_{98}={}&-x_{131}x_{136}+x_{135}x_{210}-x_{137}
\end{aligned}
\]
%END_RESIDUAL 98
%BEGIN_RESIDUAL 99
\[
\begin{aligned}
R_{99}={}&-x_{131}x_{139}+x_{138}x_{210}-x_{140}
\end{aligned}
\]
%END_RESIDUAL 99
%BEGIN_RESIDUAL 100
\[
\begin{aligned}
R_{100}={}&x_{141}-2x_{141}x_{143}+x_{144}
\end{aligned}
\]
%END_RESIDUAL 100
%BEGIN_RESIDUAL 101
\[
\begin{aligned}
R_{101}={}&-x_{141}x_{144}-x_{144}+x_{145}
\end{aligned}
\]
%END_RESIDUAL 101
%BEGIN_RESIDUAL 102
\[
\begin{aligned}
R_{102}={}&2x_{142}+x_{142}^{2}+x_{144}^{2}-x_{145}^{2}
\end{aligned}
\]
%END_RESIDUAL 102
%BEGIN_RESIDUAL 103
\[
\begin{aligned}
R_{103}={}&-x_{142}x_{147}-x_{147}+x_{148}
\end{aligned}
\]
%END_RESIDUAL 103
%BEGIN_RESIDUAL 104
\[
\begin{aligned}
R_{104}={}&-1+x_{146}^{2}+2x_{146}x_{148}+x_{147}^{2}
\end{aligned}
\]
%END_RESIDUAL 104
%BEGIN_RESIDUAL 105
\[
\begin{aligned}
R_{105}={}&x_{3}^{2}-2x_{3}x_{142}-2x_{142}+x_{149}
\end{aligned}
\]
%END_RESIDUAL 105
%BEGIN_RESIDUAL 106
\[
\begin{aligned}
R_{106}={}&-x_{3}x_{147}+x_{5}-x_{146}-x_{147}+x_{149}x_{150}
\end{aligned}
\]
%END_RESIDUAL 106
%BEGIN_RESIDUAL 107
\[
\begin{aligned}
R_{107}={}&1-2x_{4}+x_{4}^{2}-2x_{4}x_{142}+x_{151}
\end{aligned}
\]
%END_RESIDUAL 107
%BEGIN_RESIDUAL 108
\[
\begin{aligned}
R_{108}={}&-x_{4}x_{183}+x_{18}+x_{151}x_{152}-x_{182}
\end{aligned}
\]
%END_RESIDUAL 108
%BEGIN_RESIDUAL 109
\[
\begin{aligned}
R_{109}={}&-1+2x_{142}x_{194}-x_{153}+2x_{194}-x_{194}^{2}
\end{aligned}
\]
%END_RESIDUAL 109
%BEGIN_RESIDUAL 110
\[
\begin{aligned}
R_{110}={}&-1+x_{102}+x_{153}x_{154}-x_{182}-x_{183}x_{194}
\end{aligned}
\]
%END_RESIDUAL 110
%BEGIN_RESIDUAL 111
\[
\begin{aligned}
R_{111}={}&-x_{142}x_{156}-x_{156}+x_{157}
\end{aligned}
\]
%END_RESIDUAL 111
%BEGIN_RESIDUAL 112
\[
\begin{aligned}
R_{112}={}&-1+x_{155}^{2}+2x_{155}x_{157}+x_{156}^{2}
\end{aligned}
\]
%END_RESIDUAL 112
%BEGIN_RESIDUAL 113
\[
\begin{aligned}
R_{113}={}&1+x_{142}-x_{142}x_{159}-x_{159}+x_{160}
\end{aligned}
\]
%END_RESIDUAL 113
%BEGIN_RESIDUAL 114
\[
\begin{aligned}
R_{114}={}&x_{158}^{2}+2x_{158}x_{160}-2x_{159}+x_{159}^{2}
\end{aligned}
\]
%END_RESIDUAL 114
%BEGIN_RESIDUAL 115
\[
\begin{aligned}
R_{115}={}&-x_{142}x_{162}-x_{162}+x_{163}
\end{aligned}
\]
%END_RESIDUAL 115
%BEGIN_RESIDUAL 116
\[
\begin{aligned}
R_{116}={}&-1+x_{161}^{2}+2x_{161}x_{163}+x_{162}^{2}
\end{aligned}
\]
%END_RESIDUAL 116
%BEGIN_RESIDUAL 117
\[
\begin{aligned}
R_{117}={}&-x_{142}x_{165}-x_{165}+x_{166}
\end{aligned}
\]
%END_RESIDUAL 117
%BEGIN_RESIDUAL 118
\[
\begin{aligned}
R_{118}={}&-1+x_{164}^{2}+2x_{164}x_{166}+x_{165}^{2}
\end{aligned}
\]
%END_RESIDUAL 118
%BEGIN_RESIDUAL 119
\[
\begin{aligned}
R_{119}={}&-x_{142}x_{168}-x_{168}+x_{169}
\end{aligned}
\]
%END_RESIDUAL 119
%BEGIN_RESIDUAL 120
\[
\begin{aligned}
R_{120}={}&-1+x_{167}^{2}+2x_{167}x_{169}+x_{168}^{2}
\end{aligned}
\]
%END_RESIDUAL 120
%BEGIN_RESIDUAL 121
\[
\begin{aligned}
R_{121}={}&-x_{142}x_{171}-x_{171}+x_{172}
\end{aligned}
\]
%END_RESIDUAL 121
%BEGIN_RESIDUAL 122
\[
\begin{aligned}
R_{122}={}&-1+x_{170}^{2}+2x_{170}x_{172}+x_{171}^{2}
\end{aligned}
\]
%END_RESIDUAL 122
%BEGIN_RESIDUAL 123
\[
\begin{aligned}
R_{123}={}&-x_{142}x_{174}-x_{174}+x_{175}
\end{aligned}
\]
%END_RESIDUAL 123
%BEGIN_RESIDUAL 124
\[
\begin{aligned}
R_{124}={}&-1+x_{173}^{2}+2x_{173}x_{175}+x_{174}^{2}
\end{aligned}
\]
%END_RESIDUAL 124
%BEGIN_RESIDUAL 125
\[
\begin{aligned}
R_{125}={}&-x_{142}x_{177}-x_{177}+x_{178}
\end{aligned}
\]
%END_RESIDUAL 125
%BEGIN_RESIDUAL 126
\[
\begin{aligned}
R_{126}={}&-1+x_{176}^{2}+2x_{176}x_{178}+x_{177}^{2}
\end{aligned}
\]
%END_RESIDUAL 126
%BEGIN_RESIDUAL 127
\[
\begin{aligned}
R_{127}={}&-x_{142}x_{180}-x_{180}+x_{181}
\end{aligned}
\]
%END_RESIDUAL 127
%BEGIN_RESIDUAL 128
\[
\begin{aligned}
R_{128}={}&-1+x_{179}^{2}+2x_{179}x_{181}+x_{180}^{2}
\end{aligned}
\]
%END_RESIDUAL 128
%BEGIN_RESIDUAL 129
\[
\begin{aligned}
R_{129}={}&1-x_{142}x_{183}-x_{183}+x_{184}
\end{aligned}
\]
%END_RESIDUAL 129
%BEGIN_RESIDUAL 130
\[
\begin{aligned}
R_{130}={}&-1+2x_{182}+x_{182}^{2}+2x_{182}x_{184}+x_{183}^{2}
\end{aligned}
\]
%END_RESIDUAL 130
%BEGIN_RESIDUAL 131
\[
\begin{aligned}
R_{131}={}&-x_{142}x_{186}-x_{186}+x_{187}
\end{aligned}
\]
%END_RESIDUAL 131
%BEGIN_RESIDUAL 132
\[
\begin{aligned}
R_{132}={}&-1+x_{185}^{2}+2x_{185}x_{187}+x_{186}^{2}
\end{aligned}
\]
%END_RESIDUAL 132
%BEGIN_RESIDUAL 133
\[
\begin{aligned}
R_{133}={}&-x_{142}x_{189}-x_{189}+x_{190}
\end{aligned}
\]
%END_RESIDUAL 133
%BEGIN_RESIDUAL 134
\[
\begin{aligned}
R_{134}={}&-1+x_{188}^{2}+2x_{188}x_{190}+x_{189}^{2}
\end{aligned}
\]
%END_RESIDUAL 134
%BEGIN_RESIDUAL 135
\[
\begin{aligned}
R_{135}={}&-x_{142}x_{192}-x_{192}+x_{193}
\end{aligned}
\]
%END_RESIDUAL 135
%BEGIN_RESIDUAL 136
\[
\begin{aligned}
R_{136}={}&-1+x_{191}^{2}+2x_{191}x_{193}+x_{192}^{2}
\end{aligned}
\]
%END_RESIDUAL 136
%BEGIN_RESIDUAL 137
\[
\begin{aligned}
R_{137}={}&-6-2x_{0}-3x_{0}^{2}-2x_{0}x_{3}-x_{1}^{2}-2x_{3}-2x_{3}^{2}-x_{4}^{2}\\
&-2x_{5}^{2}+2x_{9}-x_{9}^{2}-x_{10}^{2}-x_{11}^{2}-x_{12}^{2}-x_{13}^{2}-x_{14}^{2}\\
&-x_{15}^{2}-x_{16}^{2}-x_{17}^{2}-x_{18}^{2}-x_{100}^{2}+2x_{102}-2x_{102}^{2}-x_{103}^{2}\\
&-x_{104}^{2}-x_{108}^{2}-x_{111}^{2}-x_{114}^{2}-x_{117}^{2}-x_{120}^{2}-x_{121}^{2}-x_{124}^{2}\\
&-x_{127}^{2}-x_{128}^{2}-x_{132}^{2}-x_{135}^{2}-x_{138}^{2}+x_{141}-x_{194}^{2}
\end{aligned}
\]
%END_RESIDUAL 137
%BEGIN_RESIDUAL 138
\[
\begin{aligned}
R_{138}={}&x_{4}-2x_{4}x_{9}-2x_{4}x_{10}-2x_{4}x_{11}-2x_{4}x_{12}-2x_{4}x_{13}-2x_{4}x_{14}-2x_{4}x_{15}\\
&-2x_{4}x_{16}-2x_{4}x_{17}+x_{194}
\end{aligned}
\]
%END_RESIDUAL 138
%BEGIN_RESIDUAL 139
\[
\begin{aligned}
R_{139}={}&-16+x_{0}-x_{195}^{2}
\end{aligned}
\]
%END_RESIDUAL 139
%BEGIN_RESIDUAL 140
\[
\begin{aligned}
R_{140}={}&-x_{142}x_{197}-x_{197}+x_{206}
\end{aligned}
\]
%END_RESIDUAL 140
%BEGIN_RESIDUAL 141
\[
\begin{aligned}
R_{141}={}&-1+x_{196}^{2}+x_{197}^{2}-x_{206}^{2}
\end{aligned}
\]
%END_RESIDUAL 141
%BEGIN_RESIDUAL 142
\[
\begin{aligned}
R_{142}={}&-x_{199}-x_{199}x_{200}+x_{207}
\end{aligned}
\]
%END_RESIDUAL 142
%BEGIN_RESIDUAL 143
\[
\begin{aligned}
R_{143}={}&-1+x_{198}^{2}+x_{199}^{2}-x_{207}^{2}
\end{aligned}
\]
%END_RESIDUAL 143
%BEGIN_RESIDUAL 144
\[
\begin{aligned}
R_{144}={}&-x_{192}^{2}+x_{208}
\end{aligned}
\]
%END_RESIDUAL 144
%BEGIN_RESIDUAL 145
\[
\begin{aligned}
R_{145}={}&x_{197}-2x_{202}x_{208}-x_{208}
\end{aligned}
\]
%END_RESIDUAL 145
%BEGIN_RESIDUAL 146
\[
\begin{aligned}
R_{146}={}&-4x_{197}-8x_{197}x_{201}+x_{200}
\end{aligned}
\]
%END_RESIDUAL 146
%BEGIN_RESIDUAL 147
\[
\begin{aligned}
R_{147}={}&-x_{142}-x_{196}x_{203}+x_{200}
\end{aligned}
\]
%END_RESIDUAL 147
%BEGIN_RESIDUAL 148
\[
\begin{aligned}
R_{148}={}&-x_{191}-x_{193}-x_{196}x_{204}+x_{198}
\end{aligned}
\]
%END_RESIDUAL 148
%BEGIN_RESIDUAL 149
\[
\begin{aligned}
R_{149}={}&-x_{100}-4x_{192}x_{205}+x_{199}
\end{aligned}
\]
%END_RESIDUAL 149
%BEGIN_RESIDUAL 150
\[
\begin{aligned}
R_{150}={}&26+2x_{0}^{2}+2x_{0}x_{3}-2x_{0}x_{147}-2x_{0}x_{156}+x_{1}^{2}-2x_{1}x_{183}+x_{3}^{2}\\
&-2x_{3}x_{156}+x_{5}^{2}-2x_{5}x_{174}+x_{9}^{2}-2x_{9}x_{159}+x_{10}^{2}-2x_{10}x_{162}+x_{11}^{2}\\
&-2x_{11}x_{165}+x_{12}^{2}-2x_{12}x_{168}+x_{13}^{2}-2x_{13}x_{171}+x_{14}^{2}-2x_{14}x_{177}+x_{15}^{2}\\
&-2x_{15}x_{180}+x_{16}^{2}-2x_{16}x_{186}+x_{17}^{2}-2x_{17}x_{189}+x_{100}^{2}-2x_{100}x_{192}-x_{142}x_{209}\\
&-2x_{146}-2x_{148}+x_{148}^{2}-2x_{155}-2x_{157}+x_{157}^{2}-2x_{158}+2x_{159}\\
&-2x_{160}+x_{160}^{2}-2x_{161}-2x_{163}+x_{163}^{2}-2x_{164}-2x_{166}+x_{166}^{2}\\
&-2x_{167}-2x_{169}+x_{169}^{2}-2x_{170}-2x_{172}+x_{172}^{2}-2x_{173}-2x_{175}\\
&+x_{175}^{2}-2x_{176}-2x_{178}+x_{178}^{2}-2x_{179}-2x_{181}+x_{181}^{2}-2x_{182}\\
&+x_{184}^{2}-2x_{185}-2x_{187}+x_{187}^{2}-2x_{188}-2x_{190}+x_{190}^{2}-2x_{191}\\
&-2x_{193}+x_{193}^{2}
\end{aligned}
\]
%END_RESIDUAL 150
%BEGIN_RESIDUAL 151
\[
\begin{aligned}
R_{151}={}&-x_{104}x_{210}+x_{128}
\end{aligned}
\]
%END_RESIDUAL 151
%BEGIN_RESIDUAL 152
\[
\begin{aligned}
R_{152}={}&-x_{121}x_{211}+x_{128}
\end{aligned}
\]
%END_RESIDUAL 152
%BEGIN_RESIDUAL 153
\[
\begin{aligned}
R_{153}={}&-2-x_{100}^{2}-x_{147}^{2}-x_{156}^{2}+2x_{159}-x_{159}^{2}-x_{162}^{2}-x_{165}^{2}\\
&-x_{168}^{2}-x_{171}^{2}-x_{174}^{2}-x_{177}^{2}-x_{180}^{2}-x_{183}^{2}-x_{186}^{2}-x_{189}^{2}\\
&+x_{192}^{2}-x_{212}^{2}-x_{213}^{2}-x_{214}-x_{214}^{2}
\end{aligned}
\]
%END_RESIDUAL 153
%BEGIN_RESIDUAL 154
\[
\begin{aligned}
R_{154}={}&-1-4x_{106}^{2}-4x_{110}^{2}-4x_{113}^{2}-4x_{116}^{2}-4x_{119}^{2}-4x_{123}^{2}-4x_{126}^{2}\\
&+x_{128}^{2}-4x_{130}^{2}-4x_{134}^{2}-4x_{137}^{2}-4x_{140}^{2}-x_{215}^{2}-x_{216}^{2}-x_{217}\\
&-x_{217}^{2}-4x_{219}^{2}
\end{aligned}
\]
%END_RESIDUAL 154
%BEGIN_RESIDUAL 155
\[
\begin{aligned}
R_{155}={}&-x_{4}x_{218}+x_{128}
\end{aligned}
\]
%END_RESIDUAL 155
%BEGIN_RESIDUAL 156
\[
\begin{aligned}
R_{156}={}&x_{2}x_{128}-x_{5}x_{218}-x_{219}
\end{aligned}
\]
%END_RESIDUAL 156
%BEGIN_RESIDUAL 157
\[
\begin{aligned}
R_{157}={}&-x_{2}x_{13}+x_{220}
\end{aligned}
\]
%END_RESIDUAL 157
%BEGIN_RESIDUAL 158
\[
\begin{aligned}
R_{158}={}&x_{102}x_{222}-2x_{103}x_{156}+2x_{108}-2x_{108}x_{159}-2x_{111}x_{162}-2x_{114}x_{165}-2x_{117}x_{168}\\
&-2x_{120}x_{171}-2x_{124}x_{177}-2x_{127}x_{180}-2x_{132}x_{186}-2x_{135}x_{189}-2x_{138}x_{192}-2x_{142}x_{222}\\
&+2x_{155}x_{156}+2x_{156}x_{157}-2x_{158}+2x_{158}x_{159}+2x_{159}x_{160}-2x_{160}+2x_{161}x_{162}+2x_{162}x_{163}\\
&+2x_{164}x_{165}+2x_{165}x_{166}+2x_{167}x_{168}+2x_{168}x_{169}+2x_{170}x_{171}+2x_{171}x_{172}+2x_{176}x_{177}\\
&+2x_{177}x_{178}+2x_{179}x_{180}+2x_{180}x_{181}+2x_{185}x_{186}+2x_{186}x_{187}+2x_{188}x_{189}+2x_{189}x_{190}\\
&+2x_{191}x_{192}+2x_{192}x_{193}+x_{221}-2x_{222}
\end{aligned}
\]
%END_RESIDUAL 158
%BEGIN_RESIDUAL 159
\[
\begin{aligned}
R_{159}={}&-1-x_{102}x_{221}+x_{103}^{2}-2x_{103}x_{155}+x_{108}^{2}-2x_{108}x_{158}+x_{111}^{2}-2x_{111}x_{161}\\
&+x_{114}^{2}-2x_{114}x_{164}+x_{117}^{2}-2x_{117}x_{167}+x_{120}^{2}-2x_{120}x_{170}+x_{124}^{2}-2x_{124}x_{176}\\
&+x_{127}^{2}-2x_{127}x_{179}+x_{132}^{2}-2x_{132}x_{185}+x_{135}^{2}-2x_{135}x_{188}+x_{138}^{2}-2x_{138}x_{191}\\
&+x_{155}^{2}-x_{156}^{2}+x_{158}^{2}+2x_{159}-x_{159}^{2}+x_{161}^{2}-x_{162}^{2}+x_{164}^{2}\\
&-x_{165}^{2}+x_{167}^{2}-x_{168}^{2}+x_{170}^{2}-x_{171}^{2}+x_{176}^{2}-x_{177}^{2}+x_{179}^{2}\\
&-x_{180}^{2}+x_{185}^{2}-x_{186}^{2}+x_{188}^{2}-x_{189}^{2}+x_{191}^{2}-x_{192}^{2}+x_{222}
\end{aligned}
\]
%END_RESIDUAL 159
%BEGIN_RESIDUAL 160
\[
\begin{aligned}
R_{160}={}&x_{102}x_{224}-2x_{104}x_{147}-2x_{121}x_{174}-2x_{128}x_{183}-2x_{142}x_{224}+2x_{146}x_{147}+2x_{147}x_{148}\\
&+2x_{173}x_{174}+2x_{174}x_{175}+2x_{182}x_{183}+2x_{183}+2x_{183}x_{184}+x_{223}-3x_{224}
\end{aligned}
\]
%END_RESIDUAL 160
%BEGIN_RESIDUAL 161
\[
\begin{aligned}
R_{161}={}&-x_{102}x_{223}+x_{104}^{2}-2x_{104}x_{146}+x_{121}^{2}-2x_{121}x_{173}+x_{128}^{2}-2x_{128}x_{182}+x_{146}^{2}\\
&-x_{147}^{2}+x_{173}^{2}-x_{174}^{2}+x_{182}^{2}-x_{183}^{2}+x_{223}+x_{224}
\end{aligned}
\]
%END_RESIDUAL 161
\endgroup


\section{Complete unit-pivot certificates}
\label{app:pivots}
For each profile, the following list contains every selected pair
$(r,v)$ of raw-system row index and raw-system variable index. The
variable order and row order of each raw system are specified in
Section~\ref{sec:definition}. Solve each selected row for its listed
variable, recursively substitute in an acyclic order, remove those rows
and variables, retain all other coordinates in increasing order, and sum
the squares of the remaining substituted residuals. There are no
existential variables for intermediate definitions. The circuit files
implement exactly this recipe, with the pooling numerators factored.
\Needspace{5\baselineskip}
\subsection*{137 integers, degree 12; paired grouping}
%BEGIN_PIVOTS integer137_degree12
\begingroup\small
\begin{longtable}{rrrrrr}
$(131,141)$ & $(15,30)$ & $(31,46)$ & $(20,35)$ & $(12,27)$ & $(27,42)$\\
$(17,32)$ & $(45,60)$ & $(19,34)$ & $(26,41)$ & $(30,45)$ & $(16,31)$\\
$(52,67)$ & $(34,49)$ & $(29,44)$ & $(8,23)$ & $(65,80)$ & $(10,25)$\\
$(23,38)$ & $(50,65)$ & $(70,85)$ & $(37,52)$ & $(40,55)$ & $(48,63)$\\
$(63,78)$ & $(6,21)$ & $(14,29)$ & $(13,28)$ & $(18,33)$ & $(56,71)$\\
$(58,73)$ & $(60,75)$ & $(68,82)$ & $(88,106)$ & $(67,84)$ & $(57,72)$\\
$(59,74)$ & $(62,77)$ & $(24,39)$ & $(1,8)$ & $(69,83)$ & $(61,76)$\\
$(82,97)$ & $(89,110)$ & $(90,113)$ & $(91,116)$ & $(92,119)$ & $(93,123)$\\
$(94,126)$ & $(98,137)$ & $(99,140)$ & $(142,192)$ & $(150,213)$ & $(81,96)$\\
$(96,130)$ & $(97,134)$ & $(140,194)$ & $(134,200)$ & $(87,107)$ & $(138,202)$\\
$(101,145)$ & $(136,201)$ & $(53,68)$ & $(151,214)$ & $(55,70)$ & $(85,99)$\\
$(83,98)$ & $(105,151)$ & $(109,157)$ & $(111,160)$ & $(113,163)$ & $(115,166)$\\
$(119,172)$ & $(121,175)$ & $(125,181)$ & $(127,184)$ & $(129,187)$ & $(103,148)$\\
$(117,169)$ & $(107,154)$ & $(95,131)$ & & &\\
\end{longtable}
\endgroup
%END_PIVOTS integer137_degree12
\Needspace{5\baselineskip}
\subsection*{121 integers, degree 16; paired grouping}
%BEGIN_PIVOTS integer121_degree16
\begingroup\small
\begin{longtable}{rrrrrr}
$(46,61)$ & $(70,85)$ & $(107,154)$ & $(143,100)$ & $(10,25)$ & $(119,172)$\\
$(92,119)$ & $(100,144)$ & $(54,69)$ & $(60,75)$ & $(42,57)$ & $(87,107)$\\
$(134,200)$ & $(8,23)$ & $(61,76)$ & $(36,51)$ & $(4,17)$ & $(68,82)$\\
$(81,96)$ & $(33,48)$ & $(149,128)$ & $(18,33)$ & $(142,192)$ & $(152,146)$\\
$(96,130)$ & $(88,106)$ & $(94,126)$ & $(16,31)$ & $(113,163)$ & $(1,8)$\\
$(19,34)$ & $(40,55)$ & $(138,202)$ & $(117,169)$ & $(97,134)$ & $(49,64)$\\
$(83,98)$ & $(51,66)$ & $(22,37)$ & $(91,116)$ & $(25,40)$ & $(56,71)$\\
$(121,175)$ & $(111,160)$ & $(73,88)$ & $(75,90)$ & $(59,74)$ & $(3,1)$\\
$(50,65)$ & $(101,145)$ & $(62,77)$ & $(109,157)$ & $(71,86)$ & $(95,131)$\\
$(44,59)$ & $(129,187)$ & $(66,81)$ & $(67,84)$ & $(123,178)$ & $(141,194)$\\
$(57,13)$ & $(82,97)$ & $(30,45)$ & $(7,22)$ & $(38,53)$ & $(28,43)$\\
$(11,26)$ & $(35,50)$ & $(127,184)$ & $(136,201)$ & $(29,44)$ & $(24,39)$\\
$(150,213)$ & $(64,79)$ & $(53,68)$ & $(105,151)$ & $(115,166)$ & $(103,148)$\\
$(63,78)$ & $(99,140)$ & $(21,36)$ & $(90,113)$ & $(77,92)$ & $(133,0)$\\
$(139,191)$ & $(93,123)$ & $(151,214)$ & $(69,83)$ & $(98,137)$ & $(125,181)$\\
$(89,110)$ & $(12,27)$ & $(85,99)$ & $(26,41)$ & $(47,62)$ & $(58,73)$\\
$(14,29)$ & & & & &\\
\end{longtable}
\endgroup
%END_PIVOTS integer121_degree16
\Needspace{5\baselineskip}
\subsection*{118 integers, degree 22; staged grouping}
%BEGIN_PIVOTS integer118_degree22
\begingroup\small
\begin{longtable}{rrrrrr}
$(151,214)$ & $(138,202)$ & $(95,131)$ & $(83,98)$ & $(81,96)$ & $(79,94)$\\
$(77,92)$ & $(75,90)$ & $(73,88)$ & $(54,69)$ & $(0,3)$ & $(145,128)$\\
$(72,87)$ & $(78,93)$ & $(74,89)$ & $(80,95)$ & $(136,201)$ & $(134,200)$\\
$(129,187)$ & $(127,184)$ & $(125,181)$ & $(121,175)$ & $(119,172)$ & $(117,169)$\\
$(115,166)$ & $(113,163)$ & $(111,160)$ & $(109,157)$ & $(105,151)$ & $(103,148)$\\
$(101,145)$ & $(87,107)$ & $(143,100)$ & $(3,1)$ & $(133,0)$ & $(150,213)$\\
$(140,194)$ & $(100,144)$ & $(99,140)$ & $(98,137)$ & $(97,134)$ & $(96,130)$\\
$(94,126)$ & $(93,123)$ & $(92,119)$ & $(91,116)$ & $(90,113)$ & $(89,110)$\\
$(71,86)$ & $(139,191)$ & $(142,192)$ & $(123,178)$ & $(85,99)$ & $(53,68)$\\
$(47,62)$ & $(44,59)$ & $(41,56)$ & $(39,54)$ & $(36,51)$ & $(28,43)$\\
$(24,39)$ & $(13,28)$ & $(4,17)$ & $(69,83)$ & $(88,106)$ & $(66,81)$\\
$(1,8)$ & $(82,97)$ & $(55,70)$ & $(61,13)$ & $(107,154)$ & $(68,82)$\\
$(64,79)$ & $(18,33)$ & $(9,24)$ & $(22,37)$ & $(11,26)$ & $(7,22)$\\
$(5,20)$ & $(63,78)$ & $(42,57)$ & $(33,48)$ & $(49,64)$ & $(37,52)$\\
$(50,65)$ & $(29,44)$ & $(25,40)$ & $(14,29)$ & $(23,38)$ & $(10,25)$\\
$(45,60)$ & $(30,45)$ & $(19,34)$ & $(17,32)$ & $(16,31)$ & $(12,27)$\\
$(20,35)$ & $(31,46)$ & $(15,30)$ & & &\\
\end{longtable}
\endgroup
%END_PIVOTS integer118_degree22
\Needspace{5\baselineskip}
\subsection*{111 integers, degree 32; staged grouping}
%BEGIN_PIVOTS integer111_degree32
\begingroup\small
\begin{longtable}{rrrrrr}
$(151,214)$ & $(138,202)$ & $(95,131)$ & $(83,98)$ & $(81,96)$ & $(79,94)$\\
$(77,92)$ & $(75,90)$ & $(73,88)$ & $(54,69)$ & $(0,3)$ & $(145,128)$\\
$(72,87)$ & $(78,93)$ & $(74,89)$ & $(80,95)$ & $(136,201)$ & $(134,200)$\\
$(129,187)$ & $(127,184)$ & $(125,181)$ & $(121,175)$ & $(119,172)$ & $(117,169)$\\
$(115,166)$ & $(113,163)$ & $(111,160)$ & $(109,157)$ & $(105,151)$ & $(103,148)$\\
$(101,145)$ & $(87,107)$ & $(143,100)$ & $(3,1)$ & $(133,0)$ & $(150,213)$\\
$(140,194)$ & $(100,144)$ & $(99,140)$ & $(98,137)$ & $(97,134)$ & $(96,130)$\\
$(94,126)$ & $(93,123)$ & $(92,119)$ & $(91,116)$ & $(90,113)$ & $(89,110)$\\
$(71,86)$ & $(139,191)$ & $(142,192)$ & $(123,178)$ & $(85,99)$ & $(53,68)$\\
$(47,62)$ & $(44,59)$ & $(41,56)$ & $(39,54)$ & $(36,51)$ & $(28,43)$\\
$(24,39)$ & $(13,28)$ & $(4,17)$ & $(69,83)$ & $(88,106)$ & $(66,81)$\\
$(1,8)$ & $(82,97)$ & $(55,70)$ & $(61,76)$ & $(107,154)$ & $(68,82)$\\
$(64,79)$ & $(18,33)$ & $(9,24)$ & $(60,75)$ & $(59,74)$ & $(57,72)$\\
$(22,37)$ & $(11,26)$ & $(7,22)$ & $(5,20)$ & $(63,78)$ & $(42,57)$\\
$(33,48)$ & $(49,64)$ & $(37,52)$ & $(48,63)$ & $(32,47)$ & $(35,50)$\\
$(50,65)$ & $(29,44)$ & $(132,188)$ & $(25,40)$ & $(51,66)$ & $(21,36)$\\
$(10,25)$ & $(6,21)$ & $(45,60)$ & $(30,45)$ & $(19,34)$ & $(26,41)$\\
$(17,32)$ & $(16,31)$ & $(12,27)$ & $(31,46)$ & &\\
\end{longtable}
\endgroup
%END_PIVOTS integer111_degree32
\Needspace{5\baselineskip}
\subsection*{106 integers, degree 38; paired grouping}
%BEGIN_PIVOTS integer106_degree38
\begingroup\small
\begin{longtable}{rrrrrr}
$(131,141)$ & $(15,30)$ & $(31,46)$ & $(20,35)$ & $(12,27)$ & $(27,42)$\\
$(17,32)$ & $(45,60)$ & $(19,34)$ & $(26,41)$ & $(30,45)$ & $(16,31)$\\
$(52,67)$ & $(34,49)$ & $(29,44)$ & $(8,23)$ & $(65,80)$ & $(10,25)$\\
$(23,38)$ & $(50,65)$ & $(70,85)$ & $(37,52)$ & $(40,55)$ & $(48,63)$\\
$(63,78)$ & $(6,21)$ & $(14,29)$ & $(25,40)$ & $(13,28)$ & $(18,33)$\\
$(56,71)$ & $(58,73)$ & $(60,75)$ & $(68,82)$ & $(88,106)$ & $(67,84)$\\
$(43,58)$ & $(57,72)$ & $(59,74)$ & $(62,77)$ & $(24,39)$ & $(7,22)$\\
$(1,8)$ & $(71,86)$ & $(69,83)$ & $(61,76)$ & $(55,70)$ & $(82,97)$\\
$(89,110)$ & $(90,113)$ & $(91,116)$ & $(92,119)$ & $(93,123)$ & $(94,126)$\\
$(98,137)$ & $(99,140)$ & $(142,192)$ & $(150,213)$ & $(66,81)$ & $(64,79)$\\
$(81,96)$ & $(96,130)$ & $(97,134)$ & $(140,194)$ & $(134,200)$ & $(87,107)$\\
$(79,94)$ & $(73,88)$ & $(75,90)$ & $(53,68)$ & $(77,92)$ & $(21,36)$\\
$(138,202)$ & $(101,145)$ & $(136,201)$ & $(33,48)$ & $(151,214)$ & $(143,193)$\\
$(85,99)$ & $(28,43)$ & $(100,144)$ & $(83,98)$ & $(36,51)$ & $(47,62)$\\
$(2,2)$ & $(105,151)$ & $(109,157)$ & $(111,160)$ & $(113,163)$ & $(115,166)$\\
$(119,172)$ & $(121,175)$ & $(125,181)$ & $(127,184)$ & $(129,187)$ & $(41,56)$\\
$(103,148)$ & $(117,169)$ & $(139,191)$ & $(107,154)$ & $(44,59)$ & $(123,178)$\\
$(133,0)$ & $(4,3)$ & $(39,54)$ & $(95,131)$ & $(11,26)$ & $(9,24)$\\
$(145,128)$ & $(152,146)$ & $(132,188)$ & $(141,142)$ & &\\
\end{longtable}
\endgroup
%END_PIVOTS integer106_degree38
\Needspace{5\baselineskip}
\subsection*{104 integers, degree 64; paired grouping}
%BEGIN_PIVOTS integer104_degree64
\begingroup\small
\begin{longtable}{rrrrrr}
$(133,0)$ & $(2,1)$ & $(0,3)$ & $(152,5)$ & $(1,8)$ & $(55,13)$\\
$(67,14)$ & $(4,17)$ & $(5,20)$ & $(6,21)$ & $(7,22)$ & $(9,24)$\\
$(10,25)$ & $(11,26)$ & $(12,27)$ & $(13,28)$ & $(14,29)$ & $(15,30)$\\
$(16,31)$ & $(17,32)$ & $(18,33)$ & $(19,34)$ & $(20,35)$ & $(21,36)$\\
$(22,37)$ & $(23,38)$ & $(25,40)$ & $(26,41)$ & $(27,42)$ & $(28,43)$\\
$(29,44)$ & $(30,45)$ & $(31,46)$ & $(32,47)$ & $(33,48)$ & $(34,49)$\\
$(35,50)$ & $(38,53)$ & $(39,54)$ & $(41,56)$ & $(42,57)$ & $(43,58)$\\
$(44,59)$ & $(46,61)$ & $(47,62)$ & $(48,63)$ & $(49,64)$ & $(50,65)$\\
$(51,66)$ & $(53,68)$ & $(56,71)$ & $(57,72)$ & $(58,73)$ & $(59,74)$\\
$(60,75)$ & $(61,76)$ & $(62,77)$ & $(63,78)$ & $(64,79)$ & $(65,80)$\\
$(66,81)$ & $(68,82)$ & $(70,85)$ & $(71,86)$ & $(72,87)$ & $(73,88)$\\
$(75,90)$ & $(77,92)$ & $(78,93)$ & $(79,94)$ & $(81,96)$ & $(82,97)$\\
$(83,98)$ & $(85,99)$ & $(86,100)$ & $(88,106)$ & $(87,107)$ & $(89,110)$\\
$(90,113)$ & $(91,116)$ & $(92,119)$ & $(93,123)$ & $(94,126)$ & $(96,127)$\\
$(145,128)$ & $(95,131)$ & $(97,132)$ & $(98,137)$ & $(99,140)$ & $(141,142)$\\
$(100,144)$ & $(101,145)$ & $(103,148)$ & $(105,151)$ & $(107,154)$ & $(109,157)$\\
$(111,160)$ & $(113,163)$ & $(115,166)$ & $(117,169)$ & $(119,172)$ & $(121,175)$\\
$(123,178)$ & $(125,181)$ & $(127,184)$ & $(142,185)$ & $(129,187)$ & $(132,188)$\\
$(139,191)$ & $(143,193)$ & $(140,194)$ & $(134,200)$ & $(138,202)$ & $(150,213)$\\
\end{longtable}
\endgroup
%END_PIVOTS integer104_degree64
\Needspace{5\baselineskip}
\subsection*{99 integers, degree 128; all5 grouping}
%BEGIN_PIVOTS integer99_degree128
\begingroup\small
\begin{longtable}{rrrrrr}
$(131,141)$ & $(15,30)$ & $(31,46)$ & $(20,35)$ & $(12,27)$ & $(27,42)$\\
$(17,32)$ & $(45,60)$ & $(19,34)$ & $(26,41)$ & $(30,45)$ & $(16,31)$\\
$(52,67)$ & $(34,49)$ & $(29,44)$ & $(8,23)$ & $(65,80)$ & $(10,25)$\\
$(23,38)$ & $(50,65)$ & $(70,85)$ & $(37,52)$ & $(40,55)$ & $(48,63)$\\
$(63,78)$ & $(6,21)$ & $(14,29)$ & $(25,40)$ & $(13,28)$ & $(18,33)$\\
$(56,71)$ & $(58,73)$ & $(60,75)$ & $(68,82)$ & $(88,106)$ & $(67,84)$\\
$(43,58)$ & $(57,72)$ & $(59,74)$ & $(62,77)$ & $(24,39)$ & $(7,22)$\\
$(1,8)$ & $(71,86)$ & $(69,83)$ & $(61,76)$ & $(55,70)$ & $(82,97)$\\
$(89,110)$ & $(90,113)$ & $(91,116)$ & $(92,119)$ & $(93,123)$ & $(94,126)$\\
$(98,137)$ & $(99,140)$ & $(142,192)$ & $(150,213)$ & $(66,81)$ & $(64,79)$\\
$(81,96)$ & $(96,130)$ & $(97,134)$ & $(140,194)$ & $(134,200)$ & $(87,107)$\\
$(79,94)$ & $(73,88)$ & $(75,90)$ & $(53,68)$ & $(77,92)$ & $(21,36)$\\
$(74,89)$ & $(78,93)$ & $(72,87)$ & $(76,91)$ & $(138,202)$ & $(101,145)$\\
$(136,201)$ & $(33,48)$ & $(151,214)$ & $(143,193)$ & $(85,99)$ & $(100,144)$\\
$(83,98)$ & $(36,51)$ & $(47,62)$ & $(2,2)$ & $(41,56)$ & $(139,191)$\\
$(44,59)$ & $(133,0)$ & $(4,17)$ & $(39,54)$ & $(95,131)$ & $(51,66)$\\
$(35,50)$ & $(32,47)$ & $(46,61)$ & $(105,151)$ & $(109,157)$ & $(111,160)$\\
$(113,163)$ & $(115,166)$ & $(119,172)$ & $(121,175)$ & $(125,181)$ & $(127,184)$\\
$(129,187)$ & $(117,169)$ & $(123,178)$ & $(107,154)$ & $(103,148)$ & $(145,128)$\\
$(132,188)$ & $(141,142)$ & $(0,3)$ & & &\\
\end{longtable}
\endgroup
%END_PIVOTS integer99_degree128
\Needspace{5\baselineskip}
\subsection*{94 integers, degree 256; all5 grouping}
%BEGIN_PIVOTS integer94_degree256
\begingroup\small
\begin{longtable}{rrrrrr}
$(131,141)$ & $(15,30)$ & $(31,46)$ & $(20,35)$ & $(12,27)$ & $(27,42)$\\
$(17,32)$ & $(45,60)$ & $(19,34)$ & $(26,41)$ & $(30,45)$ & $(16,31)$\\
$(52,67)$ & $(34,49)$ & $(29,44)$ & $(8,23)$ & $(65,80)$ & $(10,25)$\\
$(23,38)$ & $(50,65)$ & $(70,85)$ & $(37,52)$ & $(40,55)$ & $(48,63)$\\
$(63,78)$ & $(6,21)$ & $(14,29)$ & $(25,40)$ & $(13,28)$ & $(18,33)$\\
$(56,71)$ & $(58,73)$ & $(60,75)$ & $(68,82)$ & $(88,106)$ & $(67,84)$\\
$(43,58)$ & $(57,72)$ & $(59,74)$ & $(62,77)$ & $(24,39)$ & $(7,22)$\\
$(1,8)$ & $(71,86)$ & $(69,83)$ & $(61,76)$ & $(55,70)$ & $(82,97)$\\
$(89,110)$ & $(90,113)$ & $(91,116)$ & $(92,119)$ & $(93,123)$ & $(94,126)$\\
$(98,137)$ & $(99,140)$ & $(142,192)$ & $(150,213)$ & $(66,81)$ & $(64,79)$\\
$(81,96)$ & $(96,130)$ & $(97,134)$ & $(140,194)$ & $(134,200)$ & $(87,107)$\\
$(79,94)$ & $(73,88)$ & $(75,90)$ & $(53,68)$ & $(77,92)$ & $(21,36)$\\
$(74,89)$ & $(78,93)$ & $(72,87)$ & $(76,91)$ & $(138,202)$ & $(101,145)$\\
$(136,201)$ & $(33,48)$ & $(151,214)$ & $(143,193)$ & $(85,99)$ & $(28,43)$\\
$(100,144)$ & $(83,98)$ & $(36,51)$ & $(47,62)$ & $(2,2)$ & $(41,56)$\\
$(22,37)$ & $(139,191)$ & $(44,59)$ & $(133,0)$ & $(4,17)$ & $(54,69)$\\
$(39,54)$ & $(95,131)$ & $(51,66)$ & $(11,26)$ & $(9,24)$ & $(49,64)$\\
$(105,151)$ & $(109,157)$ & $(111,160)$ & $(113,163)$ & $(115,166)$ & $(119,172)$\\
$(121,175)$ & $(125,181)$ & $(127,184)$ & $(129,187)$ & $(35,50)$ & $(117,169)$\\
$(32,47)$ & $(123,178)$ & $(107,154)$ & $(103,148)$ & $(145,128)$ & $(132,188)$\\
$(141,142)$ & $(0,3)$ & & & &\\
\end{longtable}
\endgroup
%END_PIVOTS integer94_degree256
\Needspace{5\baselineskip}
\subsection*{93 integers, degree 328; staged grouping}
%BEGIN_PIVOTS integer93_degree328
\begingroup\small
\begin{longtable}{rrrrrr}
$(131,141)$ & $(15,30)$ & $(31,46)$ & $(20,35)$ & $(12,27)$ & $(27,42)$\\
$(17,32)$ & $(45,60)$ & $(19,34)$ & $(26,41)$ & $(30,45)$ & $(16,31)$\\
$(52,67)$ & $(34,49)$ & $(29,44)$ & $(8,23)$ & $(65,80)$ & $(10,25)$\\
$(23,38)$ & $(50,65)$ & $(70,85)$ & $(37,52)$ & $(40,55)$ & $(48,63)$\\
$(63,78)$ & $(6,21)$ & $(14,29)$ & $(25,40)$ & $(13,28)$ & $(18,33)$\\
$(56,71)$ & $(58,73)$ & $(60,75)$ & $(68,82)$ & $(88,106)$ & $(67,84)$\\
$(43,58)$ & $(57,72)$ & $(59,74)$ & $(62,77)$ & $(24,39)$ & $(7,22)$\\
$(1,8)$ & $(71,86)$ & $(69,83)$ & $(61,76)$ & $(55,70)$ & $(82,97)$\\
$(89,110)$ & $(90,113)$ & $(91,116)$ & $(92,119)$ & $(93,123)$ & $(94,126)$\\
$(98,137)$ & $(99,140)$ & $(142,192)$ & $(150,213)$ & $(66,81)$ & $(64,79)$\\
$(81,96)$ & $(96,130)$ & $(97,134)$ & $(140,194)$ & $(134,200)$ & $(87,107)$\\
$(79,94)$ & $(73,88)$ & $(75,90)$ & $(53,68)$ & $(77,92)$ & $(21,36)$\\
$(74,89)$ & $(78,93)$ & $(72,87)$ & $(76,91)$ & $(138,202)$ & $(101,145)$\\
$(136,201)$ & $(33,48)$ & $(151,214)$ & $(143,193)$ & $(85,99)$ & $(28,43)$\\
$(100,144)$ & $(83,98)$ & $(36,51)$ & $(47,62)$ & $(2,2)$ & $(105,151)$\\
$(109,157)$ & $(111,160)$ & $(113,163)$ & $(115,166)$ & $(119,172)$ & $(121,175)$\\
$(125,181)$ & $(127,184)$ & $(129,187)$ & $(41,56)$ & $(117,169)$ & $(22,37)$\\
$(139,191)$ & $(107,154)$ & $(44,59)$ & $(133,0)$ & $(4,17)$ & $(54,69)$\\
$(39,54)$ & $(95,131)$ & $(51,66)$ & $(123,178)$ & $(11,26)$ & $(9,24)$\\
$(49,64)$ & $(145,128)$ & $(103,148)$ & $(5,20)$ & $(3,1)$ & $(35,50)$\\
$(32,47)$ & $(132,188)$ & $(141,142)$ & $(0,3)$ & &\\
\end{longtable}
\endgroup
%END_PIVOTS integer93_degree328
\Needspace{5\baselineskip}
\subsection*{89 integers, degree 32,768; all5 grouping}
%BEGIN_PIVOTS integer89_degree32768
\begingroup\small
\begin{longtable}{rrrrrr}
$(151,214)$ & $(138,202)$ & $(95,131)$ & $(83,98)$ & $(81,96)$ & $(79,94)$\\
$(77,92)$ & $(75,90)$ & $(73,88)$ & $(54,69)$ & $(0,3)$ & $(145,128)$\\
$(72,87)$ & $(78,93)$ & $(74,89)$ & $(80,95)$ & $(76,91)$ & $(136,201)$\\
$(134,200)$ & $(129,187)$ & $(127,184)$ & $(125,181)$ & $(121,175)$ & $(119,172)$\\
$(117,169)$ & $(115,166)$ & $(113,163)$ & $(111,160)$ & $(109,157)$ & $(105,151)$\\
$(103,148)$ & $(101,145)$ & $(87,107)$ & $(143,100)$ & $(3,1)$ & $(133,0)$\\
$(150,213)$ & $(140,194)$ & $(100,144)$ & $(99,140)$ & $(98,137)$ & $(97,134)$\\
$(96,130)$ & $(94,126)$ & $(93,123)$ & $(92,119)$ & $(91,116)$ & $(90,113)$\\
$(89,110)$ & $(139,191)$ & $(71,86)$ & $(141,142)$ & $(85,99)$ & $(53,68)$\\
$(47,62)$ & $(44,59)$ & $(41,56)$ & $(39,54)$ & $(36,51)$ & $(28,43)$\\
$(24,39)$ & $(13,28)$ & $(4,17)$ & $(69,83)$ & $(88,106)$ & $(66,81)$\\
$(142,192)$ & $(123,178)$ & $(1,8)$ & $(82,97)$ & $(61,76)$ & $(67,84)$\\
$(55,70)$ & $(2,2)$ & $(64,79)$ & $(18,33)$ & $(9,24)$ & $(107,154)$\\
$(68,82)$ & $(62,77)$ & $(59,74)$ & $(60,75)$ & $(57,72)$ & $(58,73)$\\
$(56,71)$ & $(22,37)$ & $(11,26)$ & $(7,22)$ & $(63,78)$ & $(42,57)$\\
$(33,48)$ & $(5,20)$ & $(70,85)$ & $(49,64)$ & $(37,52)$ & $(48,63)$\\
$(32,47)$ & $(65,80)$ & $(46,61)$ & $(35,50)$ & $(38,53)$ & $(50,65)$\\
$(29,44)$ & $(43,58)$ & $(25,40)$ & $(51,66)$ & $(14,29)$ & $(132,188)$\\
$(21,36)$ & $(6,21)$ & $(10,25)$ & $(8,23)$ & $(23,38)$ & $(40,55)$\\
$(34,49)$ & $(27,42)$ & $(52,67)$ & $(30,45)$ & $(19,34)$ & $(17,32)$\\
$(16,31)$ & $(26,41)$ & $(12,27)$ & $(20,35)$ & $(31,46)$ & $(15,30)$\\
$(131,141)$ & & & & &\\
\end{longtable}
\endgroup
%END_PIVOTS integer89_degree32768
\Needspace{5\baselineskip}
\subsection*{88 integers, degree 65,540; all5 grouping}
%BEGIN_PIVOTS integer88_degree65540
\begingroup\small
\begin{longtable}{rrrrrr}
$(131,141)$ & $(15,30)$ & $(31,46)$ & $(20,35)$ & $(12,27)$ & $(27,42)$\\
$(17,32)$ & $(45,60)$ & $(19,34)$ & $(26,41)$ & $(30,45)$ & $(16,31)$\\
$(52,67)$ & $(34,49)$ & $(29,44)$ & $(8,23)$ & $(65,80)$ & $(10,25)$\\
$(23,38)$ & $(50,65)$ & $(70,85)$ & $(37,52)$ & $(40,55)$ & $(48,63)$\\
$(63,78)$ & $(6,21)$ & $(14,29)$ & $(25,40)$ & $(13,28)$ & $(18,33)$\\
$(56,71)$ & $(58,73)$ & $(60,75)$ & $(68,82)$ & $(88,106)$ & $(67,84)$\\
$(43,58)$ & $(57,72)$ & $(59,74)$ & $(62,77)$ & $(24,39)$ & $(7,22)$\\
$(1,8)$ & $(71,86)$ & $(69,83)$ & $(61,76)$ & $(55,70)$ & $(82,97)$\\
$(89,110)$ & $(90,113)$ & $(91,116)$ & $(92,119)$ & $(93,123)$ & $(94,126)$\\
$(98,137)$ & $(99,140)$ & $(142,192)$ & $(150,213)$ & $(66,81)$ & $(64,79)$\\
$(81,96)$ & $(96,130)$ & $(97,134)$ & $(140,194)$ & $(134,200)$ & $(87,107)$\\
$(79,94)$ & $(73,88)$ & $(75,90)$ & $(53,68)$ & $(77,92)$ & $(21,36)$\\
$(74,89)$ & $(78,93)$ & $(72,87)$ & $(76,91)$ & $(138,202)$ & $(101,145)$\\
$(136,201)$ & $(33,48)$ & $(151,214)$ & $(143,193)$ & $(85,99)$ & $(28,43)$\\
$(100,144)$ & $(83,98)$ & $(36,51)$ & $(47,62)$ & $(2,2)$ & $(41,56)$\\
$(22,37)$ & $(139,191)$ & $(44,59)$ & $(133,0)$ & $(4,17)$ & $(54,69)$\\
$(39,54)$ & $(95,131)$ & $(51,66)$ & $(11,26)$ & $(9,24)$ & $(49,64)$\\
$(5,20)$ & $(3,1)$ & $(80,95)$ & $(105,151)$ & $(109,157)$ & $(111,160)$\\
$(113,163)$ & $(115,166)$ & $(119,172)$ & $(121,175)$ & $(125,181)$ & $(127,184)$\\
$(129,187)$ & $(35,50)$ & $(117,169)$ & $(32,47)$ & $(123,178)$ & $(107,154)$\\
$(103,148)$ & $(86,100)$ & $(145,128)$ & $(42,57)$ & $(132,188)$ & $(141,142)$\\
$(38,53)$ & $(0,3)$ & & & &\\
\end{longtable}
\endgroup
%END_PIVOTS integer88_degree65540
\Needspace{5\baselineskip}
\subsection*{87 integers, degree 8,388,612; all5 grouping}
%BEGIN_PIVOTS integer87_degree8388612
\begingroup\small
\begin{longtable}{rrrrrr}
$(0,3)$ & $(1,8)$ & $(2,2)$ & $(3,1)$ & $(4,17)$ & $(5,20)$\\
$(6,21)$ & $(7,22)$ & $(8,23)$ & $(9,24)$ & $(10,25)$ & $(11,26)$\\
$(12,27)$ & $(13,28)$ & $(14,29)$ & $(15,30)$ & $(16,31)$ & $(17,32)$\\
$(18,33)$ & $(19,34)$ & $(20,35)$ & $(21,36)$ & $(22,37)$ & $(23,38)$\\
$(24,39)$ & $(25,40)$ & $(26,41)$ & $(27,42)$ & $(28,43)$ & $(29,44)$\\
$(30,45)$ & $(31,46)$ & $(32,47)$ & $(33,48)$ & $(34,49)$ & $(35,50)$\\
$(36,51)$ & $(37,52)$ & $(38,53)$ & $(39,54)$ & $(40,55)$ & $(41,56)$\\
$(42,57)$ & $(43,58)$ & $(44,59)$ & $(45,60)$ & $(46,61)$ & $(47,62)$\\
$(48,63)$ & $(49,64)$ & $(50,65)$ & $(51,66)$ & $(52,67)$ & $(53,68)$\\
$(54,69)$ & $(55,70)$ & $(56,71)$ & $(57,72)$ & $(58,73)$ & $(59,74)$\\
$(60,75)$ & $(61,76)$ & $(62,77)$ & $(63,78)$ & $(64,79)$ & $(65,80)$\\
$(66,81)$ & $(67,84)$ & $(68,82)$ & $(69,83)$ & $(70,85)$ & $(71,86)$\\
$(72,87)$ & $(73,88)$ & $(74,89)$ & $(75,90)$ & $(76,91)$ & $(77,92)$\\
$(78,93)$ & $(79,94)$ & $(80,95)$ & $(81,96)$ & $(82,97)$ & $(83,98)$\\
$(85,99)$ & $(86,100)$ & $(87,107)$ & $(88,106)$ & $(89,110)$ & $(90,113)$\\
$(91,116)$ & $(92,119)$ & $(93,123)$ & $(94,126)$ & $(95,131)$ & $(96,130)$\\
$(97,134)$ & $(98,137)$ & $(99,140)$ & $(100,144)$ & $(101,145)$ & $(103,148)$\\
$(105,151)$ & $(107,154)$ & $(109,157)$ & $(111,160)$ & $(113,163)$ & $(115,166)$\\
$(117,169)$ & $(119,172)$ & $(121,175)$ & $(123,178)$ & $(125,181)$ & $(127,184)$\\
$(129,187)$ & $(131,141)$ & $(132,188)$ & $(133,0)$ & $(134,200)$ & $(136,201)$\\
$(138,202)$ & $(139,191)$ & $(140,194)$ & $(141,142)$ & $(142,192)$ & $(143,193)$\\
$(145,128)$ & $(150,213)$ & $(151,214)$ & & &\\
\end{longtable}
\endgroup
%END_PIVOTS integer87_degree8388612


\section{The compact companion's four fixed coefficients}
\label{app:coefficients}
The direct integer family does not use this appendix. For the compact
positive-domain family, the following finite recipe supplies the actual
four coefficients, including the exact seed. The accompanying independent
constructor reproduces the stored bytes and their hashes.
% Exact finite recipe extracted from the retained indexed_constants.tex.
% In this appendix K denotes the fixed port K_0, and all other symbols are constants.
\begingroup
\providecommand{\fl}[1]{\left\lfloor#1\right\rfloor}
\providecommand{\rg}[2]{\{#1,\ldots,#2\}}
In this appendix $K$ denotes the port $K_0$ in the new equations.
Every symbol introduced here is a fixed constant, independent of the ten
quantified variables. The formulas and 436-character seed are the retained
indexed BFS943 radix67 specification; they reconstruct the exact supplied
payloads, with $\beta=2^{590069}-1$.

\textbf{The four fixed integers.}
Every index, set, sum, product and floor below defines a constant, not
another unknown. Put
\[
\rho=2^{67},\quad B=\rho^{8807},\quad\beta=B-1,\quad
G_m(x)=\sum_{j=0}^{m-1}x^j.
\]
Define
\[
\begin{split}
\mathcal E={}&\rg0{942}\cup\rg{983}{992}\cup\rg{1003}{1011}\\
&\cup\rg{1013}{1031}\cup\{1063,1064,1095,2195,3296\}.
\end{split}
\]
Then
\begin{align*}
M_C={}&B-1-\sum_{e\in\mathcal E\setminus\{1\}}\rho^e-2\rho^{1063},\\
M_F={}&B-1+4+(2^{51}-2)(\rho^{4393}+\rho^{5495})\\
&+\rho^{1013}G_{10}(\rho)+\rho^{1065}G_9(\rho)
  +\rho^{2195}+\rho^{3296}\\
&+(2^{65}-1)\rho^{8792}G_{13}(\rho)+(2^{19}-1)\rho^{8805},\\
K={}&\rho^{30}+\rho^{62}+\rho^{1100}+\rho^{2201}+\rho^{8849}\\
&+(\rho^{4393}+\rho^{5495})\Phi(1/\rho),
\end{align*}
where
\[
\begin{split}
\Phi(x)={}&(1+2^{51})G_{943}(x)
 -2^{10}\sum_{i=0}^{942}v(t_i)x^i\\
&+\sum_{(b,e,m)\in\mathcal C}2^b x^e G_m(2x),\\
\mathcal C={}&\{(10,1023,8),(18,1013,9),(27,1003,8),\\
&(35,983,9),(44,1031,1),(45,1022,1),(46,1011,1),\\
&(47,992,1),(48,1095,1),(49,2195,1),(50,3296,1)\}.
\end{split}
\]
The rational intermediate in K cancels exactly; all four results are positive integers.

\textbf{The 943 indexed fixed tiles.}
Use the actions $u_j,v_j,d_j$ decoded below. Let
\[
\begin{gathered}
h(q,b)=2(q+1)+b,\quad
\mathcal R=\{u_j:d_j=1,\ 0\le j<430\},\\
\mathcal L=\{u_j:d_j=0,\ 0\le j<430\},\\
r_q=6+|\mathcal R\cap\rg0{q-1}|,\quad
l_q=205+|\mathcal L\cap\rg0{q-1}|.
\end{gathered}
\]
An interval ending at $-1$ is empty. The table defines $t_i$ for
$0\le i\le942$; all index ranges are fixed:
\[
\setlength{\arraycolsep}{1.4pt}\begin{array}{c|l}
\text{index }i&\text{tile }t_i=(w,n,e,s)\\\hline
0,1&(2,435,3,1),\ (5,0,0,0)\\
j+2\ (j=2q+b<430)&((1-d_j)l_{u_j},h(q,b),d_jr_{u_j},v_j)\\
432+b\ (b=0,1)&(0,h(215,b),0,h(215,b))\\
434+b\ (b=0,1)&(0,b,0,b)\\
424+2r_q+b\ (q\in\mathcal R,b=0,1)&(r_q,b,0,h(q,b))\\
424+2l_q+b\ (q\in\mathcal L,b=0,1)&(0,b,l_q,h(q,b))\\
924+j\ (0\le j<4)&(0,j+430\fl{j/2},0,438)\\
928&(2,434,2,0)\\
929&(3,436,4,4)\\
930&(4,436,4,0)\\
497+t\ (t=434,435,436)&(1,438,1,t)\\
\substack{934+3a+b-6\fl{a/4}\\a\in\{0,1,4\},\ b=0,1,2}&(a,437,b,437)
\end{array}
\]
For $t=(w,n,e,s)$ put
\[
v(t)=w+2^8n+2^{17}e+2^{25}s+127\,2^{34}.
\]

\textbf{A closed decoder for the 430 actions.}
Put $D_b(A,k)=\fl{A/b^k}\bmod b$, where
$a\bmod b=a-b\fl{a/b}$. For $0\le q<215$, $b=0,1$, $j=2q+b$:
\[
\begin{gathered}
N=\mathcal A\bmod2^{636},\quad
T_0=\fl{\mathcal A/2^{636}}\bmod2^{1675},\\
E_0=\fl{\mathcal A/2^{2311}},\\
\nu_j=D_2(N,j),\quad \sigma_j=\sum_{k=0}^{j-1}\nu_k,\\
m_q=D_2(N,430+q),\quad \tau_q=\sum_{k=0}^{q-1}m_k,\\
F_q=14(1-m_q)+m_qD_{16}(E_0,\tau_q),\\
u_j=\nu_j(1+\sigma_j)+(1-\nu_j)D_{216}(T_0,j-\sigma_j),\\
v_j=D_2(F_q,2b),\quad d_j=D_2(F_q,2b+1).
\end{gathered}
\]
Empty sums are zero. Here $(v_0,d_0,u_0)=(1,1,1)$; the constructor tile
uses this actual first target, including its state renaming.

\textbf{The integer $\mathcal A$.}
The alphabet, in order of digit values $0,\ldots,61$, is:\par
{\fontsize{9}{10}\selectfont\ttfamily 0123456789ABCDEFGHIJKLMNOPQRSTUVWXYZabcdefghijklmnopqrstuvwxyz}.\par
Concatenate the following 436 characters as digits $e_0,\ldots,e_{435}$;
put $\mathcal A=\sum_{j=0}^{435}e_j62^{435-j}$.\par
\begingroup\fontsize{9}{10}\selectfont\ttfamily
BxEnhvy6NipzkWmROo0GsA3uK3twWLxscLq20UTumy5Cdu461LwD6g7RDNzF8Y\par
iYDFTpB8ehAypd7d6OsOL4APamnJeNXQypUezCmW9n4oXP6XfgmiWuAWZ8ZJ7U\par
OP7gu4kVdlGZqZHULQMar1dXNeJW4i7oNVI8usp8BlovjjZbwPCovZI328fRUq\par
4q92R81s5GjWQaQfGhkq9UYMK3rRMaBmsOIBkJqc2cj2VRIZ6mjzkBIOkW9SjZ\par
FzeA6J5u4j3sBL3SdgwXzM9mHvQivu8EJ1TBTpRhHDqZIbmYgwQ8VqTAFotl6F\par
We1YC554GKrVHsA8oUGKSA6F6k5KREBSD0Tgd7bFALY84LLbpw574BQyvdoHrl\par
0Oe35uwDf15DpT6VJWZBjSbHPpPAINYTte0eUxnTFTKBlGj8Am6sb1gLldh8WZ\par
8P\par
\endgroup

\endgroup


\end{document}
